\pdfinfoomitdate=1
\pdftrailerid{}
\pdfsuppressptexinfo=15
\documentclass[11pt,letterpaper]{article}
\usepackage[T1]{fontenc}
\usepackage{lmodern}
\usepackage{amsmath,amssymb,amsthm,mathtools}
\usepackage[margin=1in]{geometry}
\usepackage{booktabs,longtable,array}
\usepackage{microtype}
\usepackage[colorlinks=true,linkcolor=black,citecolor=black,urlcolor=blue]{hyperref}
\usepackage{xurl}
\setlength{\parskip}{4pt}
\setlength{\parindent}{1.2em}
\setlength{\emergencystretch}{2em}
\setcounter{tocdepth}{1}
\numberwithin{equation}{section}
\newtheorem{theorem}{Theorem}[section]
\newtheorem{lemma}[theorem]{Lemma}
\newtheorem{proposition}[theorem]{Proposition}
\newtheorem{corollary}[theorem]{Corollary}
\theoremstyle{definition}\newtheorem{definition}[theorem]{Definition}
\theoremstyle{remark}\newtheorem{remark}[theorem]{Remark}
% Macros of the two sources; both define every one of them identically.
\newcommand{\PP}{\mathbb P}
\newcommand{\EE}{\mathbb E}
\newcommand{\one}{\mathbf1}
\newcommand{\eps}{\varepsilon}
\newcommand{\ceil}[1]{\left\lceil #1\right\rceil}
\newcommand{\floor}[1]{\left\lfloor #1\right\rfloor}
\DeclareMathOperator{\Aut}{Aut}
\DeclareMathOperator{\Poi}{Poisson}
% Added in the merge (batch 106).
\newcommand{\file}[1]{\nolinkurl{#1}}
\newenvironment{writenote}[1][]{\par\smallskip\begingroup\small
 \noindent\textbf{[write]}\if\relax\detokenize{#1}\relax\else~\textbf{#1.}\fi\ }%
 {\par\endgroup\smallskip}
\newcommand{\wtag}{\textup{\textbf{[write]}}}
\newenvironment{partabstract}{\begin{quote}\small\noindent\textbf{Abstract of this Part (as delivered).}\ }{\end{quote}}
\newcommand{\partsource}[1]{\par\noindent\begin{minipage}{\textwidth}\small #1\end{minipage}\par\medskip}
\hypersetup{pdftitle={Unlabeled circle graphs (OEIS A156808 and A156809)},pdfauthor={},pdfsubject={Merged report a156808-circle-graphs (batch 106): the leading equivalent with error bounds and inverses (Report 150), the first corrections and refined inverses (Report 152)}}
\title{\textbf{Unlabeled circle graphs\\ (OEIS A156808 and A156809)}\\[7pt]
\large The leading equivalent with error bounds and inverses,\\ and the first corrections}
\author{Two research reports of one external research session\\(bundle Reports 150 and 152), merged into one ProveIt report}
\date{Parts I and II: 3 October 2026\quad merged: 5 October 2026}
\begin{document}
\hypersetup{pageanchor=false}
\maketitle
\begin{abstract}
Let $c_n$ and $g_n$ count connected and all unlabeled circle graphs on $n$ vertices (OEIS A156808 and A156809), and put $h_n=(2n-1)!!/(4n)$. This report prints two manuscripts as Parts~I and~II, the second a sequel to the first. Part~I proves $c_n\sim g_n\sim e^{-3/2}h_n$, matching an upper bound from reciprocal representation fibres with an injective construction from asymmetric split-prime cores carrying leaf, true-twin and false-twin decorations. It also proves the connected fraction $1-1/(2n)+O(n^{-2})$, relative errors $O(1/n)$ for both counts through a self-contained Poisson approximation of local matching patterns, independent limiting $\Poi(1/2)$ decoration counts, a $\Poi(3/2)$ core deficit, the asymmetry probability $e^{-1}$, a labeled equivalent, and a Lambert-$W$ inverse with ceiling-safe brackets; and it derives $\mathrm{A156808}(13)=21{,}593{,}488{,}017$ from the OEIS value of $\mathrm{A156809}(13)$. Its prime probability $e^{-3}+o(1)$ follows from lemmas of Bassino, Bouvel, F\'eray, Gerin and Pierrot, who state the bound $e^{-4}$; Part~I says that this is not a new theorem. Part~II proves $c_n=e^{-3/2}h_n(1-\frac{15}{4n}+O(n^{-2}))$, $g_n=e^{-3/2}h_n(1-\frac{13}{4n}+O(n^{-2}))$ and $c_n/g_n=1-\frac1{2n}-\frac1{n^2}+O(n^{-3})$, from a uniform local expansion that gives the prime probability $e^{-3}(1-5/n+O(n^{-2}))$ and a canonical split-tree transfer to graph classes. It adds first corrections for automorphism-weighted and labeled counts and for rare automorphism groups, a refined inverse, and a transfer theorem for every fixed order that is conditional on the corresponding expansion of the prime counts; its second-order coefficients depend on an unknown prime constant $\beta$. The value of $\beta$, an unconditional expansion beyond the first correction and effective constants remain open. The report is unrefereed, and nothing in it is formalized.
\end{abstract}
\hypersetup{pageanchor=true}
{\small\setlength{\parskip}{0pt}\tableofcontents}
\newpage

\section*{Guide to this report}\phantomsection\label{cgr:sec:guide}
\addcontentsline{toc}{section}{Guide to this report}
The two Parts are two manuscripts of the session bundle of Reports 1--243 (arrival commit \texttt{60f54ea06}), placed in this directory by commit \texttt{47fc7a069} (batch~106). They study one pair of sequences, the numbers $c_n$ (OEIS A156808) and $g_n$ (OEIS A156809) of connected and of all unlabeled circle graphs of order $n$, in one model: the indexed perfect matchings of $2n$ cyclically ordered endpoints, $m_n=(2n-1)!!$ of them, acted on by the dihedral group of order $4n$, with the normalization $h_n=m_n/(4n)$.
\begin{center}
\small
\begin{tabular}{@{}l l >{\raggedright\arraybackslash}p{0.47\textwidth} l@{}}
\toprule
Part & Bundle report & Delivered title & Labels\\
\midrule
I & 150 (base) & Leading asymptotics of unlabeled circle graphs & \texttt{cgr:ld:}\\
II & 152 & First corrections for unlabeled circle graphs & \texttt{cgr:fc:}\\
\bottomrule
\end{tabular}
\end{center}
\emph{Arrangement.} Report~150 is the base of the merge (its \file{Report150.tex} was staged as this \file{article.tex}) and is printed first, as Part~I: it introduces the setting, proves the leading equivalent and its relative $O(1/n)$ error, and alone proves the typical-core laws, the explicit local total-variation bound and the derived value $c_{13}$. Report~152 is a sequel, not a new edition: it says that ``the earlier Report~150 established the leading equivalent and relative $O(1/n)$ bounds; this report identifies the individual first coefficients'' (Section~\ref{cgr:fc:sub:priority}). It is printed second, as Part~II. Its archive carries no copy of Report~150. Dependency order is also the order of writing.

\emph{Numbering.} Sections are numbered continuously; statements and equations are numbered within sections, as in both manuscripts. Part~I keeps Report~150's Sections~1--12, statement and equation numbers; its Appendix~A is Section~\ref{cgr:ld:app:dependency}. In Part~II a section number is Report~152's plus~$13$ (Sections~14--24), and so are the first components of its statement and equation numbers: its Theorems~1.1 and~1.2 are Theorems~\ref{cgr:fc:thm:main} and~\ref{cgr:fc:thm:transfer-summary}, its Theorem~3.1 is Theorem~\ref{cgr:fc:thm:local}, its Lemma~5.1 is Lemma~\ref{cgr:fc:lem:three}, its Theorem~6.1 is Theorem~\ref{cgr:fc:thm:finite-transfer}, its Theorem~10.1 is Theorem~\ref{cgr:fc:thm:inverse}, and its equation $(k.j)$ is equation $((k+13).j)$ here. Remark~\ref{cgr:ld:rem:symmetry-credit}, the last numbered statement of its section, and Section~\ref{cgr:sec:further}, which closes the report, were added in the write; no delivered number moved.

Notes added when the manuscripts were merged are set in small type and tagged \textbf{[write]}, with their date (5~October 2026). They give cross-references between the Parts, mark statements that Part~II supersedes or answers, credit prior work, and record what the write checked. Superseded statements and answered questions are printed as delivered, with a dated note; no delivered statement was changed and no delivered symbol was renamed. Part~I's open questions are its Section~\ref{cgr:ld:sec:limits}; Part~II states its limits in its last Section~\ref{cgr:fc:sec:checks}; Section~\ref{cgr:sec:further} gathers both under Vladimir's standing rule of 4~October 2026, with their sources, sketches and what is missing.

\section*{What is proved, and where}\phantomsection\label{cgr:sec:status}
\addcontentsline{toc}{section}{What is proved, and where}
Here $m_n=(2n-1)!!$, $h_n=m_n/(4n)$, $p_n^{\mathrm a}$ counts asymmetric split-prime circle graphs of order $n$, and $u=L/W_0(2L/e)$ with $L$ the logarithm of the threshold.
\begin{center}
\small
\begin{tabular}{@{}>{\raggedright\arraybackslash}p{0.55\textwidth} >{\raggedright\arraybackslash}p{0.40\textwidth}@{}}
\toprule
Statement & Where\\
\midrule
$c_n\sim g_n\sim e^{-3/2}h_n$; $g_n-c_n=g_{n-1}+O(m_n/n^3)$; $c_n/g_n=1-\frac1{2n}+o(n^{-1})$ & Part~I, Theorem~\ref{cgr:ld:thm:main} (Sections~\ref{cgr:ld:sec:symmetry}--\ref{cgr:ld:sec:disconnected})\\
Prime probability $e^{-3}+o(1)$ (the source states $>e^{-4}$) & Part~I, \eqref{cgr:ld:eq:prime-prob}, from \cite[Lemmas 5.7--5.10]{BBFGP24}\\
$S_n\le4n(2en)^{n/2}e^{\sqrt n}$ for the symmetric matchings & Part~I, Lemma~\ref{cgr:ld:lem:symmetry}; restated in Part~II, \eqref{cgr:fc:eq:symmetry}; method of \cite[Proposition 5.14]{BBFGP24} (\wtag\ Remark~\ref{cgr:ld:rem:symmetry-credit})\\
Typical recovered core: decoration counts $\to\Poi(1/2)^{\otimes3}$, deficit $\to\Poi(3/2)$; $\Aut(G)\cong(C_2)^{t+f}$; asymmetric, leafless, split-prime with limiting probabilities $e^{-1}$, $e^{-1/2}$, $e^{-3/2}$; labeled counts $\sim e^{-2}(2n)!/(2^{n+2}n)$ & Part~I, Theorem~\ref{cgr:ld:thm:core}, Proposition~\ref{cgr:ld:prop:aut}, \eqref{cgr:ld:eq:other-prob}, \eqref{cgr:ld:eq:labeled}\\
Relative errors $O(1/n)$ for $c_n$, $g_n$, the laws above and the labeled counts; $c_n/g_n=1-\frac1{2n}+O(n^{-2})$; $d_{\rm TV}((X_n,Y_n,Z_n),\Poi(1)^{\otimes3})\le\min\{1,200/n\}$ for $n\ge6$ & Part~I, Theorem~\ref{cgr:ld:thm:quantitative}, Lemma~\ref{cgr:ld:lem:stein}, \eqref{cgr:ld:eq:explicit-TV}\\
Smooth inverse $\phi(x)=u+1+\frac{3+\log2}{2\log(2u)}+O(\frac1{u\log u})$; least-index brackets of width $O(1/(u\log u))$ inside two ceilings & Part~I, Propositions~\ref{cgr:ld:prop:inverse}, \ref{cgr:ld:prop:fixed-bracket}, \eqref{cgr:ld:eq:quant-inverse}\\
$c_{13}=21{,}593{,}488{,}017$, derived by Euler inversion & Part~I, Section~\ref{cgr:ld:sec:checks}\\
$\EE[(1+x)^X(1+y)^Y(1+z)^Z]=e^{x+y+z}(1+F/M+O(M^{-2}))$ uniformly on compacts; $\PP(X=Y=Z=0)=e^{-3}(1-10/M+O(M^{-2}))$ & Part~II, Theorem~\ref{cgr:fc:thm:local}\\
Three-chord interval lemma; $\PP(\text{indecomposable})=e^{-3}(1-5/n+O(n^{-2}))$; $p_n^{\mathrm a}=e^{-3}h_n(1-5/n+O(n^{-2}))$ & Part~II, Lemma~\ref{cgr:fc:lem:three}, \eqref{cgr:fc:eq:prime-prob}, \eqref{cgr:fc:eq:prime-main}\\
Finite canonical split-tree transfer; rooted counts $b_0,\ldots,b_3=1,3,11,58$ & Part~II, Theorem~\ref{cgr:fc:thm:finite-transfer}, Subsection~\ref{cgr:fc:sub:rooted}\\
$c_n=e^{-3/2}h_n(1-\frac{15}{4n}+O(n^{-2}))$, $g_n=e^{-3/2}h_n(1-\frac{13}{4n}+O(n^{-2}))$; in the Stirling normalization $-\frac{91}{24}$, $-\frac{79}{24}$ & Part~II, Theorem~\ref{cgr:fc:thm:main}\\
$c_n/g_n=1-\frac1{2n}-\frac1{n^2}+O(n^{-3})$; the transfer to every fixed order, \emph{conditional} on the prime expansion to that order; conditional second coefficients $\beta-\frac{277}{32}$, $\beta-\frac{297}{32}$ & Part~II, Theorem~\ref{cgr:fc:thm:transfer-summary}, Section~\ref{cgr:fc:sec:allorders}\\
Automorphism-weighted counts $c_n(t)$, $g_n(t)$; labeled counts with $-\frac{47}{12n}$, $-\frac{41}{12n}$; $\PP(\Aut(G)=1)=e^{-1}(1-\frac1{2n}+O(n^{-2}))$; $\EE|\Aut(G)|^{-t}$; $\PP(\Aut(G)\text{ nonabelian})=\frac1{2n}+O(n^{-2})$ & Part~II, Section~\ref{cgr:fc:sec:weighted}\\
Refined inverse with centres $s_a$, $K_c=\frac{103}{24}$, $K_g=\frac{91}{24}$, width $O(1/(u^2\log u))$ & Part~II, Theorem~\ref{cgr:fc:thm:inverse}\\
\midrule
\multicolumn{2}{@{}l}{\emph{Open in both Parts}}\\
\multicolumn{2}{@{}p{0.96\textwidth}@{}}{The prime second coefficient $\beta$, hence the second graph coefficients; an unconditional expansion to every fixed order; effective constants, onsets and certified thresholds; reconciliation with the 2026 journal text of the main source; an audit of the labeled consequences; unbounded automorphism observables. See Section~\ref{cgr:sec:further}.}\\
\bottomrule
\end{tabular}
\end{center}
The two companions check finite identities, enumerations, normalizations and coefficient algebra in exact arithmetic; no asymptotic statement rests on them. The leading inverses of both Parts are instances of the transseries volume's factorial core; their two-ceiling brackets are not instances of its staircase theorem (next section).

\section*{The OEIS entries, and the inverses}\phantomsection\label{cgr:sec:oeis}
\addcontentsline{toc}{section}{The OEIS entries, and the inverses}
The entries as read by the write on 5~October 2026 are A156808 revision \#9 (June 29, 2026) and A156809 revision \#16 (June 30, 2026), the revisions of Part~I's snapshot (shipped as \file{data/150-leading-companion-sources-oeis_snapshot.json}).
\begin{itemize}
\item \textbf{A156808} (Lars Eirik Danielsen, Feb 16 2009), ``Number of nonisomorphic connected circle graphs of order $n$'': terms for $n=1,\ldots,12$, ending with $892278076$; no formula, comment or asymptotic statement. Part~I derives $c_{13}=21{,}593{,}488{,}017$; the value is not in the entry, and nothing was submitted to the OEIS.
\item \textbf{A156809} (Lars Eirik Danielsen, Feb 16 2009), ``Number of nonisomorphic circle graphs of order $n$'': terms for $n=1,\ldots,13$; the extension line credits $a(13)=22576188846$ to Tom Johnston, Oct 17 2020, whose blog post ``Enumerating circle graphs'' (2020, read by the write) gives that count for all circle graphs of order~13 and no connected count. No formula, comment or asymptotic statement.
\end{itemize}
Both entries cite Danielsen and Parker \cite{DP09}, whose Table~3 gives the counts through order~12. The theorems of this report are additions to the entries, not corrections of them.

\emph{The inverses.} Both Parts write the threshold through $u$, the solution of $u(\log(2u)-1)=L$ (Part~I's \eqref{cgr:ld:eq:core-root}, Part~II's \eqref{cgr:fc:eq:inverse-vars}). This is the factorial core \texttt{p0:prop:factorial-core} of the transseries volume \cite{TSvol}, $X(\kappa\log X+d)=L$, with $\kappa=1$ and $d=\log2-1$; its solution $\exp(W_0(Le^{d})-d)$ equals $L/W_0(2L/e)=u$, and its side condition $\log(2u)>1$ holds for large $L$. After taking logarithms it is the Lambert core \texttt{p0:thm:lambert-core} with $a=b=1$, in the variable $\log(2u)-1=W_0(2L/e)$ and with right-hand side $\log L+\log2-1$. The corrections to $u$ (Part~I's $1+\frac{3+\log2}{2\log(2u)}$, Part~II's further $(K_a-B^2/(2D^2))/(uD)$) are derived in the Parts by Taylor expansion of the explicit Gamma-function model $H(t)$ at $u$; they are not taken from the volume, and we do not cite them as instances of its centred expansions. The two-ceiling brackets \eqref{cgr:ld:eq:eta-bracket}, \eqref{cgr:ld:eq:quant-inverse} and \eqref{cgr:fc:eq:inverse-bracket} follow the separation pattern of \texttt{p0:thm:staircase}\,(2) but are not instances of it: that theorem inverts an admissible interpolation of the sequence itself, while the Parts compare $a_n$ with the model $H$, which does not interpolate $a_n$ ($H(n)=e^{-3/2}h_n\ne a_n$), and prove the brackets directly at the integers. Neither manuscript cites the volume.

\section*{Notation across the two Parts}\phantomsection\label{cgr:sec:notation}
\addcontentsline{toc}{section}{Notation across the two Parts}
Common to both Parts, with the same meaning: $c_n$, $g_n$ (offsets one; $g_0=1$, $m_0=1$); $m_n=(2n-1)!!$; $h_n=m_n/(4n)$; $\mathcal M_n$, $G_D$, $R_n(G)$; $S_n$, the number of indexed matchings fixed by a nonidentity dihedral symmetry; $p_n^{\mathrm a}$ (Part~I: asymmetric split-prime circle graphs; Part~II: ``abstractly asymmetric prime circle graphs'', the same class); $M=2n$ in the local analysis; $X,Y,Z$ (and Part~I's $X_n,Y_n,Z_n$) the counts of chords at cyclic distance one, chords at distance two, and pairs of consecutive endpoints matched to another consecutive pair; the smooth model $H(t)=e^{-3/2}2^t\Gamma(t+\frac12)/(4\sqrt\pi\,t)$; $u=L/W_0(2L/e)$ and $N_a$, the least index. The symbols below change meaning between the Parts or inside one; each Part's text fixes its meaning there, and each Part opens with a short table of its own readings. No symbol was renamed.
\begingroup\small
\renewcommand{\theHtable}{cgr-notation.\arabic{table}}% unique anchor for the uncaptioned longtable
\begin{longtable}{@{}>{\raggedright\arraybackslash}p{0.13\textwidth} >{\raggedright\arraybackslash}p{0.83\textwidth}@{}}
\toprule
Symbol & Meanings\\
\midrule
\endhead
$c_r$, $c_n$ & $c_n=\mathrm{A156808}(n)$ throughout. Part~II, Section~\ref{cgr:fc:sec:polymer}: $c_r$ is the coefficient of $u^r$ in the partition polynomial $H_M(u)$ \eqref{cgr:fc:eq:partition}. \emph{Tempting false reading:} that $c_r$ is not a circle-graph count. Part~II also writes $c_n(t)$ for the automorphism-weighted count.\\
$B$ & Part~I: $B=\frac32(1+\log2)$ in \eqref{cgr:ld:eq:logH}; $B_0,\ldots,B_3$ the sums of Section~\ref{cgr:ld:sec:quantitative}. Part~II: $B=\sum_{P\sim Q}w_Pw_Q$, $B_1(s)$, $B_*$ (Section~\ref{cgr:fc:sec:hardcore}); $B_d(z)$ the rooted-branch polynomial (Section~\ref{cgr:fc:sec:split}); $B_1=-a^2-a/2+2$ (Subsection~\ref{cgr:fc:sub:second}); $B=\frac{3+\log2}2$ in \eqref{cgr:fc:eq:inverse-vars}. \emph{False reading:} Part~II's inverse constant $B$ is Part~I's $B-\log2$, not Part~I's $B$.\\
$A$, $D$ & Part~I: $A=\log(2u)$ (Section~\ref{cgr:ld:sec:inverse}); $D$ a matching, $D_0$, $D^\alpha$; $\mathcal D_{n;\ell,t,f}$ the decorated families. Part~II: $D=\log(2u)$ (Part~I's $A$); $D_n$ a uniform matching; $\mathcal D=s\,d/ds$; $A(s)$ the leading polymer sum, $A=Pe^{3/2}$, a $\Poi(3/2)$ variable $A$, $A_1,A_2$ polynomials, $A_x=N_{S_x}(\rho_x)$.\\
$q$, $\rho$ & Part~I: $q_s=m_sm_{n-s}/m_n$ \eqref{cgr:ld:eq:q-def}; $\rho_r=\prod_{j<r}(M-2j-1)^{-1}$ \eqref{cgr:ld:eq:pattern-prob}. Part~II: $q_r=\prod_{j<r}(M-2j-1)^{-1}$ \eqref{cgr:fc:eq:matching-exact}, which is Part~I's $\rho_r$, not its $q_s$; $q$ an expansion order (Section~\ref{cgr:fc:sec:allorders}); $\rho=2^{-t}$ and $\rho_x$ a root vertex (Sections~\ref{cgr:fc:sec:split}, \ref{cgr:fc:sec:weighted}).\\
$K$ & Part~I: a truncation level, fixed in Section~\ref{cgr:ld:sec:upper} and $\ceil{\log n}$ in Section~\ref{cgr:ld:sec:quantitative}. Part~II: $K_c=\frac{103}{24}$, $K_g=\frac{91}{24}$ in the inverse centres.\\
$L$, $T$, $F$ & Part~I: the decoration types $L,T,F$, their counts $L_n,T_n,F_n$; $L=\log x$; $F(t)=t(\log(2t)-1)$ in the proof of Proposition~\ref{cgr:ld:prop:inverse}. Part~II: $F(x,y,z)$ the local polynomial \eqref{cgr:fc:eq:F}; $F_\delta(t)=\log H(t)+\delta/t$; $L=\log y$; $L_1,L_2$ polynomials; $T$ a split tree.\\
$H$ & Part~I: $H$ a prime core graph (Section~\ref{cgr:ld:sec:lower}); $H(t)$ the smooth model \eqref{cgr:ld:eq:H}. Part~II: $H_M(u)$ the partition polynomial; $H$ a prime centre (Section~\ref{cgr:fc:sec:split}); the same $H(t)$.\\
$\alpha$, $\beta$, $\gamma$, $\eta$ & Part~I: $\alpha,\beta$ index local patterns (Section~\ref{cgr:ld:sec:quantitative}); $\eta>0$ and $\eta(x)$ bracket errors (Section~\ref{cgr:ld:sec:inverse}). Part~II: $\alpha$ and $\beta$ the first and second relative prime coefficients ($\alpha=-5$ proved, $\beta$ unknown), $\alpha_t$; $\gamma=\alpha+5/4$, $\eta=\gamma+1/2$.\\
$C$ & A constant in \eqref{cgr:ld:eq:uniform-g} and in \eqref{cgr:fc:eq:coarse} (the same role); Part~I: $C_1,C_2,C_*$ inverse constants. Part~II: $C$ also a constant in \eqref{cgr:fc:eq:q-global} and in the polymer bounds, $C_*$ a remainder sum (Section~\ref{cgr:fc:sec:hardcore}), $C_a,y_a$ inverse constants; $C_n^{\mathrm{lab}}$, $G_n^{\mathrm{lab}}$ the labeled counts, which Part~I writes $c_n^{\mathrm{lab}}$, $g_n^{\mathrm{lab}}$.\\
$S$ & $S_n$ as above (both). Part~I: $S_\alpha$ a pattern support. Part~II: $S$ a bound $|u|\le S/M$, $S_P$ a polymer support, $(S_x,\rho_x)$ a rooted branch, $S_3$ the symmetric group.\\
$x$, $y$, $z$ & Part~I: $x$ the threshold input. Part~II: $x,y,z$ complex weights (Section~\ref{cgr:fc:sec:polymer}); $y$ the threshold input (Section~\ref{cgr:fc:sec:inverse}).\\
$d$, $k$, $r$, $s$ & Part~I: $d$ a rotation order and a root displacement; $k$ a decomposition size or decoration count; $r$ a deletion size or arc parameter, $r_n$ a log error; $s=r+1$, $s(x)$ the inverse centre. Part~II: $d$ a rotation order, a branch-size bound, $d_n^k$ the decomposed-matching count, $d_0,d_1$ displacements; $r$ a split-side bound, $r(A)$ a chord count, $r_s$ rooted classes, $r_\pm$ roots; $s$ a complex variable, a split side, $s_a(y)$ the inverse centres.\\
$P$, $W$ & Part~I: $P\sim\Poi(1)$ (Proposition~\ref{cgr:ld:prop:aut}); $W_0$ Lambert. Part~II: $P$ the prime amplitude ($e^{-3}$), $P,Q$ polymers; $W$ the polymer sum, $W_1(s)$, $W_*$; $W_0$ Lambert.\\
$b$ & Part~II only: $b_j$ the rooted connected circle-graph counts, $b_j(t)$ their weighted versions; $b=a_2$ an excess-two multiplicity (Subsection~\ref{cgr:fc:sub:second} and Section~\ref{cgr:fc:sec:coeff}).\\
\bottomrule
\end{longtable}
\addtocounter{table}{-1}% the uncaptioned longtable must not consume a table number
\endgroup

\section*{Provenance and merge decisions}\phantomsection\label{cgr:sec:provenance}
\addcontentsline{toc}{section}{Provenance and merge decisions}
Two manuscripts of the session bundle of Reports 1--243. Their author lines read ``Report 150'' and ``Report 152''; neither names an author, a tool or an addressee (each sets an empty PDF author field), neither carries ``prepared for private review'' wording, and neither pins a ProveIt commit or names a repository path. They arrived in commit \texttt{60f54ea06} and were placed by \texttt{47fc7a069} (batch~106).
\begin{itemize}
\item \textbf{Part~I}, bundle Report~150, archive \file{Circle_Graph_Enumeration_Error_Bounds_and_Inverses_Source.zip} (622,124 bytes, 18 files; \file{Report150.tex}, 1445 lines, 23 pages, dated 3 October 2026); the base. Contributes the setting, the symmetry and deletion bounds, the fibre identity, the leaf lemmas, the upper bound by independent leaf moves, the lower bound by canonical recovery of decorations, the transfer to disconnected graphs, saturation with the core and automorphism laws and the labeled equivalent, the forcing coupling with the self-contained Poisson approximation and the quantitative sandwich, the smooth inverse with its brackets, and the finite checks with the derived $c_{13}$.
\item \textbf{Part~II}, bundle Report~152, archive \file{Circle_Graph_First_Corrections_and_Refined_Inverses_Source.zip} (630,050 bytes, 18 files; \file{Report152.tex}, 1200 lines, 19 pages, dated 3 October 2026). Contributes the uniform local expansion through the polymer representation, the global matching-probability remainder, the three-chord interval lemma and the first prime correction, the finite canonical split-tree transfer with its tails, the first graph coefficients, the second connectivity term, the conditional transfer to every fixed order and the conditional second coefficients, the automorphism-weighted and labeled corrections, and the refined inverse.
\end{itemize}
Where the merge had to choose:
\begin{enumerate}
\item \emph{Order and base.} Base and dependency order agree: Report~150 is Part~I, its sequel Report~152 Part~II (see the Guide).
\item \emph{Restatements.} Part~II's Section~\ref{cgr:fc:sec:inputs} restates in its own words the structural inputs Part~I takes from \cite{BBFGP24} and \cite{KZ15}, including the faithfulness argument for prime chord actions, and re-proves Part~I's symmetry bound (Lemma~\ref{cgr:ld:lem:symmetry}) and coarse bound \eqref{cgr:ld:eq:coarse-g}--\eqref{cgr:ld:eq:uniform-g}, so that its proof does not depend on Part~I. The two manuscripts share about 1.5\% (of Report~150) and 1.8\% (of Report~152) of their word 8-grams (intake measurement). The restatements are printed in full, with notes pointing to Part~I.
\item \emph{Supersession.} Part~II's Theorem~\ref{cgr:fc:thm:main} identifies the first relative coefficients that Part~I's Theorem~\ref{cgr:ld:thm:quantitative} bounds; its Theorem~\ref{cgr:fc:thm:transfer-summary} sharpens \eqref{cgr:ld:eq:connected} and \eqref{cgr:ld:eq:quant-main}; its Theorem~\ref{cgr:fc:thm:inverse} refines Part~I's Section~\ref{cgr:ld:sec:inverse}; its Section~\ref{cgr:fc:sec:weighted} refines Part~I's asymmetry probability and labeled equivalent. Part~I's superseded statements stay, with dated notes. Only Part~I has the $\Poi(1/2)^{\otimes3}$ laws with total-variation rates, the explicit bound $200/n$, the leafless and split-prime probabilities, saturation, and $c_{13}$.
\item \emph{Front matter and bibliography.} The two title blocks and abstracts are printed at the heads of their Parts, with each manuscript's scope paragraph; one table of contents replaces two. The two bibliographies are merged: the entries for \cite{BBFGP24}, \cite{KZ15} and the two OEIS entries occur in both manuscripts with different annotations, and both annotations are kept; \cite{DP09} is Part~I's, \cite{GP12} Part~II's; the write added \cite{TSvol}.
\item \emph{Appendix.} Report~150's Appendix~A (``Dependency map and reproducible release'') is printed as Section~\ref{cgr:ld:app:dependency}, so that Part~II's sections follow it.
\item \emph{Write additions.} This front matter, the reading-conventions tables of the Parts, Remark~\ref{cgr:ld:rem:symmetry-credit}, Section~\ref{cgr:sec:further}, the dated notes, labels for unlabelled sections, and the label prefixes. Everything else is delivered text.
\end{enumerate}

\section*{What was checked, and what was not}\phantomsection\label{cgr:sec:trust}
\addcontentsline{toc}{section}{What was checked, and what was not}
The report is unrefereed and not formalized. At intake (dossier 106-GRAPHS of batch~106, 5~October 2026) both manuscripts were read in full and no step was found false. Checked there: the Euler transform of A156808 gives A156809 through order~12; the derived $c_{13}$; the constant $B=\frac32(1+\log2)$ from Stirling's formula for $\Gamma(t+\frac12)$; that the upper and lower constants of Part~I agree; $F(-1,-1,-1)=-10$; $\EE[-A^2+3A/2]=-\frac32$ for $A\sim\Poi(\frac32)$ and the arithmetic giving $-\frac{15}4$, $-\frac{13}4$, $K_c$, $K_g$, $-\frac{91}{24}$, the labeled values $b_1(1)=2$, $b_2(1)=\frac{19}3$ and $\frac{13}{12}$; and the delivered suites, whose recorded outputs were regenerated byte for byte on copies (README). The intake also ran a Monte Carlo simulation of $\PP(X=Y=Z=0)$ (note at Theorem~\ref{cgr:fc:thm:local}; uncertified).

The write read the cited statements of the external inputs: in arXiv:2402.06394v1 of Bassino, Bouvel, F\'eray, Gerin and Pierrot (the file whose SHA-256 both manuscripts record; arXiv lists no later version as of 5~October 2026) Definition~5.3, Proposition~5.5, Lemmas~5.6--5.10, Propositions~5.11, 5.12 and~5.14 and the bound~(12); in Klav\'ik and Zeman's STACS paper (the recorded file) Lemma~7 ($\Aut$ of a split tree equals $\Aut$ of its graph) and the proof sketch of Lemma~8 on p.~549, which says that a prime circle graph has a unique circle representation up to rotations and reflections, so that its automorphism group is a subgroup of the dihedral group; in Gioan and Paul's arXiv:0810.1823v2 Lemmas~2.3 and~2.8, Corollary~2.6, Theorem~2.9 and Corollary~2.10 (the published version was not compared); and Table~3 of Danielsen and Parker's arXiv:0804.2576v2. Each says what the Parts take from it. The write recomputed $c_{13}$ and the ratios quoted in the notes to Section~\ref{cgr:ld:sec:checks}; by brute force over all graphs on two to five vertices, the rooted counts $b_0,\ldots,b_3$ and the root-stabilizer orders used in Section~\ref{cgr:fc:sec:weighted}; the sum $353.96<354$ behind \eqref{cgr:ld:eq:explicit-TV}; and, numerically, the transfer arithmetic of \eqref{cgr:fc:eq:secondc} (note there). Remark~\ref{cgr:ld:rem:symmetry-credit} is proved by the write.

What was \emph{not} done: the analytic arguments (the Poisson approximation of Lemma~\ref{cgr:ld:lem:stein}, the polymer bounds and the hard-core remainder, the global matching remainder, the split-tree transfer and its tails, the inverse expansions) were read and their displayed algebra checked, but they are the manuscripts' proofs, corroborated by finite checks, not independently verified. Not read: the 2026 journal version of Bassino et al.\ (\emph{Electronic Journal of Probability} \textbf{31}, article~110), at intake or at the write (its publisher page refused automated access); Gabor, Supowit and Hsu's paper, cited through \cite[Proposition 5.12]{BBFGP24}; and, beyond its abstract, the paper of Arratia, Bollob\'as, Coppersmith and Sorkin. Priority is therefore bounded as both manuscripts bound it: their results were not located in the sources checked, which is not a claim of novelty against the final journal text.

\section*{Relation to neighbouring reports and formal projects}\phantomsection\label{cgr:sec:neighbours}
\addcontentsline{toc}{section}{Relation to neighbouring reports and formal projects}
No other report of the collection treats circle graphs or A156808 and A156809 (repository search, 5~October 2026). Two reports placed in the same commit and written separately in batch~106 are siblings: \file{oeis-sequence-asymptotics/a123448-permutation-graphs} (bundle Report~153, unlabeled permutation graphs, OEIS A123448) and \file{oeis-sequence-asymptotics/a005975-interval-graphs} (bundle Report~133, unlabeled interval graphs, OEIS A005975 and A005976). The three share a pipeline: count a representation model (here indexed matchings; there permutations and interval orders), show that almost every graph has a recoverable prime core whose representations form one full orbit of the representation symmetry group (here dihedral of order $4n$; there the Klein four-group and a duality of order two), divide by that group's order, make the symmetric cores negligible, and invert with ceiling-safe Lambert-$W$ brackets. They share no statement: the decomposition theories (split trees here, modular decomposition there), the external inputs and the normalizations ($(2n-1)!!/(4n)$, $n!/4$, half the Fishburn numbers) differ, no manuscript cites another, and a merged report would give several symbols two to four meanings (placement commit \texttt{47fc7a069}). Bassino et al.\ \cite{BBFGP24} treat permutation graphs as well. The multigraph reports \file{a007716-bipartite-multigraphs} and \file{a307316-leafless-multigraphs} use the same rare-symmetry philosophy for other objects and share no statement with this one. The leading inverses are instances of the transseries volume's factorial core (above). Placement in the collection confers no formal status: no Lean or Rocq file in the repository treats circle graphs, and no statement of this report is formalized.
\clearpage
\part{Leading asymptotics of unlabeled circle graphs}\label{cgr:ld:part}
\partsource{Bundle Report~150, archive \file{Circle_Graph_Enumeration_Error_Bounds_and_Inverses_Source.zip}, delivered as \file{Report150.tex} (1445 lines, 23 pages), the base of the merge. Delivered title block ``Leading asymptotics of unlabeled circle graphs / Connected and all graphs in OEIS A156808 and A156809 / Report 150 / 3 October 2026''. Printed as delivered, with \textbf{[write]} notes; section, statement and equation numbers are those of the delivery, and its Appendix~A is Section~\ref{cgr:ld:app:dependency}. Remark~\ref{cgr:ld:rem:symmetry-credit} is added in the write.}
\begin{center}
\small
\begin{tabular}{@{}>{\raggedright\arraybackslash}p{0.46\textwidth} >{\raggedright\arraybackslash}p{0.46\textwidth}@{}}
\toprule
Reading conventions of Part~I & Elsewhere in the report\\
\midrule
$B=\frac32(1+\log2)$ in \eqref{cgr:ld:eq:logH}; $B_0,\ldots,B_3$ the sums of Section~\ref{cgr:ld:sec:quantitative} & Part~II: $B=\frac{3+\log2}2$ (this Part's $B-\log2$) in the inverse; $B$ also a conflict sum and $B_d(z)$ a branch polynomial\\
$A=\log(2u)$ (Section~\ref{cgr:ld:sec:inverse}) & Part~II: $D=\log(2u)$; $A$ an amplitude, a Poisson variable, polynomials\\
$q_s=m_sm_{n-s}/m_n$; $\rho_r$ the probability of $r$ prescribed chords & Part~II: $q_r$ is this Part's $\rho_r$\\
$K$ a truncation level & Part~II: $K_c$, $K_g$ inverse constants\\
$L,T,F$ decoration types; $L=\log x$; $H$ a core graph and $H(t)$ the model & Part~II: $F(x,y,z)$, $F_\delta$; $L=\log y$; $H_M(u)$\\
$\alpha,\beta$ pattern indices; $\eta$, $\eta(x)$ bracket errors & Part~II: $\alpha=-5$, $\beta$ prime coefficients; $\eta=\gamma+1/2$\\
$c_n^{\mathrm{lab}}$, $g_n^{\mathrm{lab}}$ & Part~II: $C_n^{\mathrm{lab}}$, $G_n^{\mathrm{lab}}$\\
\bottomrule
\end{tabular}
\end{center}
\begin{partabstract}
Let $c_n$ and $g_n$ count connected and arbitrary unlabeled circle
graphs on $n$ vertices. We prove
\[
 c_n\sim g_n\sim e^{-3/2}\frac{(2n-1)!!}{4n}.
\]
The proof combines published local Poisson limits and uniqueness of
prime representations with two model-specific arguments: inverse
representation-fiber weights give an upper bound, while canonically
recoverable leaf and twin decorations of asymmetric prime cores give
the matching lower bound. We also prove the connected fraction
$1-1/(2n)+o(1/n)$, three independent limiting Poisson decoration
counts of mean $1/2$, a Poisson core deficit of mean $3/2$, and
asymmetry probability $e^{-1}$. A smooth Lambert-$W$ inverse is
converted into ceiling-safe least-index brackets. A self-contained forcing-coupling argument sharpens both individual
count estimates to relative error $O(n^{-1})$ and the connectivity
error to $O(n^{-2})$. Fresh exact finite checks and a deterministic
source package accompany the proofs.
\end{partabstract}
\noindent\textbf{Scope.} The complete primary texts inspected were
arXiv:2402.06394v1 (2024), an additional complete author-hosted 2024
copy, and the references specified below. The 2026 journal version
was identified bibliographically but its final text was not inspected.
The result fills a leading-equivalent gap in the sources checked;
this is not an exhaustive priority claim. The proof gives neither an
all-orders expansion nor an effective finite-$n$ onset. Finite checks
are not substitutes for the structural proof. No external submission,
author contact, or public repository change is part of this report.
\begin{writenote}[Scope, 5~October 2026]
The bounded review stated here is kept. The write read arXiv:2402.06394v1 at the places this Part cites; arXiv lists no later version as of 5~October 2026. The journal text was read neither at intake nor at the write (``What was checked, and what was not'' in the front matter). The first correction coefficients that this Part leaves open are identified in Part~II.
\end{writenote}

\section{The result and its counting conventions}\label{cgr:ld:sec:result}
A circle graph is the intersection graph of finitely many chords of a
circle, with two vertices adjacent exactly when their corresponding
chords cross. All graphs in this report are finite, simple and
undirected. Chord endpoints can be assumed distinct. An unlabeled
graph means an isomorphism class, counted once regardless of its
automorphism group.

Let $\mathcal M_n$ be the perfect matchings of a fixed cyclically
ordered endpoint set $\{0,1,\ldots,2n-1\}$. Its cardinality is
\begin{equation}\label{cgr:ld:eq:mn}
 m_n=(2n-1)!!=\frac{(2n)!}{2^n n!},\qquad m_0=1.
\end{equation}
These are \emph{indexed} matchings: the endpoints are indexed, but
the chords do not have independent labels. Write $G_D$ for the
unlabeled intersection graph of $D\in\mathcal M_n$. The dihedral
group acting on the endpoints has order $4n$. A matching is
\emph{geometrically asymmetric} when its stabilizer in that group is
trivial. A graph is \emph{asymmetric} when its graph automorphism
group is trivial. These two notions must not be interchanged for
general circle graphs.

Set
\begin{equation}\label{cgr:ld:eq:sequences}
 c_n=\mathrm{A156808}(n),\quad
 g_n=\mathrm{A156809}(n),\quad
 h_n=\frac{m_n}{4n}\quad(n\ge1).
\end{equation}
Both OEIS sequences begin at $n=1$. We use $g_0=1$ for the empty
graph when discussing disjoint unions.

\begin{theorem}[Leading enumeration]\label{cgr:ld:thm:main}
As $n\to\infty$,
\begin{equation}\label{cgr:ld:eq:main}
 c_n\sim g_n\sim e^{-3/2}h_n
 \sim\frac{e^{-3/2}}{2\sqrt2\,n}\left(\frac{2n}{e}\right)^n.
\end{equation}
Moreover,
\begin{align}
 g_n-c_n&=g_{n-1}+O(m_n/n^3)
          =c_{n-1}+O(m_n/n^3),\label{cgr:ld:eq:gap}\\
 \frac{c_n}{g_n}&=1-\frac1{2n}+o(n^{-1}).\label{cgr:ld:eq:connected}
\end{align}
Section~\ref{cgr:ld:sec:quantitative} strengthens the individual equivalents
to relative error $O(n^{-1})$ and the error in
\eqref{cgr:ld:eq:connected} to $O(n^{-2})$. The individual first
correction coefficients are not identified.
\end{theorem}
\begin{writenote}[Answered in Part~II, 5~October 2026]
Part~II's Theorem~\ref{cgr:fc:thm:main} identifies the individual first corrections: $c_n=e^{-3/2}h_n(1-\frac{15}{4n}+O(n^{-2}))$ and $g_n=e^{-3/2}h_n(1-\frac{13}{4n}+O(n^{-2}))$. Its Theorem~\ref{cgr:fc:thm:transfer-summary} sharpens \eqref{cgr:ld:eq:connected} to $1-\frac1{2n}-\frac1{n^2}+O(n^{-3})$. This theorem and Section~\ref{cgr:ld:sec:quantitative} are printed as delivered.
\end{writenote}

The main asymptotic normalization is $h_n$, rather than $m_n$.
Dividing all matchings by $4n$ alone does not count circle graphs:
a graph may have several inequivalent chord representations, and
geometric symmetries can shorten individual orbits. Sections
\ref{cgr:ld:sec:symmetry}--\ref{cgr:ld:sec:lower} handle both issues explicitly.

\section{Published inputs and the bounded source review}\label{cgr:ld:sec:inputs}
\subsection{Splits and prime representations}
A split of a graph is a partition $V=A\sqcup B$ whose crossing
edges form $A'\times B'$ for some $A'\subseteq A$ and
$B'\subseteq B$; the crossing set is allowed to be empty. The split
is nontrivial if $|A|,|B|\ge2$. A graph is \emph{split-prime} if it
has no nontrivial split.

Bassino, Bouvel, F\'eray, Gerin and Pierrot
\cite[Definition 5.3 and Proposition 5.5]{BBFGP24} define matching
decompositions by four successive cyclic intervals, allowing empty
intervals, such that every chord lies within one interval or joins
opposite intervals. A $k$-decomposition uses the source's anchored convention: $C_1$
contains endpoint $0$ (called $1$ in the source), the other
nonempty intervals follow in cyclic order, and $k$ counts the
chords in $C_2\cup C_4$. Thus $k$ and $n-k$ are not freely
interchangeable in the cited probability bound. A decomposition is
nontrivial when $2\le k\le n-2$. Their equivalence is
\begin{equation}\label{cgr:ld:eq:prime-equivalence}
 D\text{ is indecomposable}
 \quad\Longleftrightarrow\quad G_D\text{ is split-prime}.
\end{equation}
We accept this published structural theorem. We also use their
Proposition 5.12, attributed there to Gabor, Supowit and Hsu:
a split-prime circle graph of order at least five has one matching
representation up to rotation and reflection.

For the automorphism statement we use the natural geometric action,
not just an isomorphism between abstract groups. Klav\'ik and Zeman
\cite[proof of Lemma 8 and the following geometric explanation,
pp.~548--549]{KZ15} describe prime graph automorphisms as rotations
or reflections of the unique circle representation. Thus each prime
graph automorphism is induced by such a symmetry.

For completeness, no nonidentity geometric symmetry can fix every
chord of a split-prime representation of order at least five. A
rotation fixing every chord would have order two and all chords
would be diameters, giving a complete graph. A reflection fixing
every chord has all but at most one chord perpendicular to its
axis; these chords do not cross one another. The possible chord on
the axis gives at most a star together with isolated vertices. None
of these graphs is split-prime at this order. The natural action on
chords is therefore faithful. In particular,
\begin{equation}\label{cgr:ld:eq:prime-asym}
 \text{a prime matching is geometrically asymmetric}
 \quad\Longleftrightarrow\quad
 \text{its graph is asymmetric}.
\end{equation}
No such equivalence is asserted outside the prime class.

\subsection{The local Poisson input}\label{cgr:ld:sub:poisson}
For a uniform $D_n\in\mathcal M_n$, let $X_n$ count chords joining
consecutive endpoints and $Y_n$ count chords joining endpoints at
cyclic distance two. Thus a chord counted by $Y_n$ has exactly one
endpoint in its shorter open arc, for all sufficiently large $n$.
Let $Z_n$ count unordered pairs of disjoint consecutive endpoint
pairs whose endpoints are matched to one another; each such local
configuration is counted once, with either the crossing or
noncrossing pairing allowed. These definitions agree with
\cite[Section 5.2.4]{BBFGP24} in the asymptotic range.

The precise published inputs are
\begin{align}
 (X_n,Y_n,Z_n)&\xrightarrow{d}(X,Y,Z),\label{cgr:ld:eq:poi-source}\\
 X,Y,Z&\text{ are independent }\Poi(1),\notag\\
 \PP(D_n\text{ has a }k\text{-decomposition for some }k>2)
   &=O(n^{-1}),\label{cgr:ld:eq:large-decomp}\\
 \PP(D_n\text{ has no }2\text{-decomposition})&\longrightarrow e^{-3}.
 \label{cgr:ld:eq:no-two}
\end{align}
These are respectively Lemma 5.10 (proved in Appendix A.2),
Lemma 5.7, and Lemma 5.8 of \cite{BBFGP24}. Their Lemma 5.9
identifies the simultaneous vanishing of $X_n,Y_n,Z_n$ with the
absence of the two endpoint decomposition sizes. The union bound
and \eqref{cgr:ld:eq:large-decomp}--\eqref{cgr:ld:eq:no-two} immediately give
\begin{equation}\label{cgr:ld:eq:prime-prob}
 \PP(D_n\text{ is indecomposable})=e^{-3}+o(1).
\end{equation}
Although their Proposition 5.11 states only the eventual lower
bound $e^{-4}$, this sharper limit follows directly from the
preceding lemmas. It is not a new Poisson theorem proved here.
\begin{writenote}[What the source states, 5~October 2026]
Checked in arXiv:2402.06394v1. Its Proposition~5.11 states that, for $n$ large enough, the probability that a uniform matching of size $n$ is indecomposable exceeds $e^{-4}$. Its proof bounds the probability of decomposability by that of a $2$-decomposition plus that of a $k$-decomposition with $k>2$, and Lemmas~5.7 and~5.8 make this sum tend to $1-e^{-3}$. So the limit $e^{-3}$ in \eqref{cgr:ld:eq:prime-prob} follows from the source's own lemmas (the upper bound from Lemma~5.8 alone), as this Part says: it sharpens a stated bound and corrects nothing in the source. The source uses $e^{-4}$ for its lower bound~(12), that the number of circle graphs of order $n$ is at least $e^{-4}m_n/(4n)$ for large $n$; Theorem~\ref{cgr:ld:thm:main} replaces that bound by an equivalent. Part~II sharpens \eqref{cgr:ld:eq:prime-prob} to $e^{-3}(1-5/n+O(n^{-2}))$ in \eqref{cgr:fc:eq:prime-prob}.
\end{writenote}

\subsection{What was and was not inspected}\label{cgr:ld:sub:inspected}
The version actually used is arXiv:2402.06394v1, dated 9 February
2024. A complete author-hosted copy with PDF metadata dated
13 February 2024 was also checked; its relevant text agrees with
that version and still gives the coarse bound. Journal metadata
lists \emph{Electronic Journal of Probability} \textbf{31} (2026),
article 110, DOI \href{https://doi.org/10.1214/26-EJP1570}
{10.1214/26-EJP1570}. The final journal text was not obtained and
is not being cited as inspected.

Danielsen and Parker \cite[Section 3, Table 3]{DP09} supply exact
circle-graph counts through order twelve. The inspected version is
arXiv:0804.2576v2 (1 October 2009). The OEIS entries
\cite{OEISc,OEISg} provide the same connected data and one additional
all-graph value at order thirteen. Neither inspected OEIS entry
states the equivalent \eqref{cgr:ld:eq:main}. An older OEIS-linked paper
of Arratia, Bollob\'as, Coppersmith and Sorkin was checked only at
abstract level; its Euler-circuit enumeration is not used as a
theorem about graph-isomorphism classes.

The literature claim is deliberately bounded: the equivalent was
not located in these inspected texts or the accompanying targeted
catalogue search. This is neither proof of universal novelty nor a
claim to settle a well-known open problem. A final-version source
reconciliation remains useful future work.
\begin{writenote}[The sources as read at the write, 5~October 2026]
Table~3 of arXiv:0804.2576v2 (p.~8) agrees with Table~\ref{cgr:ld:tab:counts} through order~12 and with the shipped extraction \file{data/150-leading-companion-sources-danielsen_parker_table3.txt}, which also carries the table's two columns of local-complementation orbit counts. The live OEIS entries are the revisions of the shipped snapshot (A156808 \#9, A156809 \#16); the value $g_{13}=22576188846$ is credited there to Tom Johnston (Oct 17 2020), whose blog post ``Enumerating circle graphs'' (2020) gives the all-graph count at order~13 only. Neither entry has a formula or a comment. The ``coarse bound'' above is the source's Proposition~5.11, the bound $e^{-4}$ (shipped provenance record \file{data/150-leading-SOURCE_PROVENANCE.json}), which the source turns into its lower bound~(12) on the number of circle graphs (note after \eqref{cgr:ld:eq:prime-prob}). Remark~\ref{cgr:ld:rem:symmetry-credit} relates the source's Proposition~5.14 to Lemma~\ref{cgr:ld:lem:symmetry}.
\end{writenote}

\section{Symmetry bounds and the exact fiber identity}\label{cgr:ld:sec:symmetry}
\subsection{Every geometric exception is negligible}
Let $S_n$ be the number of indexed matchings fixed by some
nonidentity symmetry of the regular $2n$-gon.

\begin{lemma}\label{cgr:ld:lem:symmetry}
For every fixed $A>0$, $S_n=o(m_n n^{-A})$.
In particular, $nS_n/m_n\to0$.
\end{lemma}
\begin{proof}
A rotation of order $d\ge2$ has $k=2n/d$ endpoint cycles. If a
matching is invariant, it pairs two distinct cycles using one of
$d$ offsets, or pairs a cycle internally. The latter is possible
only when $d$ is even and must pair opposite points. Hence its
number of fixed matchings is
\begin{equation}\label{cgr:ld:eq:rotation-egf}
 k![z^k]\exp\left(\one_{2\mid d}z+\frac d2z^2\right).
\end{equation}
For a power series with nonnegative coefficients,
$[z^k]F(z)\le F(r)r^{-k}$. Choosing $r=\sqrt{k/d}$ and using
$k!\le k^k$, the quantity in \eqref{cgr:ld:eq:rotation-egf} is at most
\begin{align}
 (kd)^{k/2}\exp(k/2+\sqrt{k/d})
 &\le (2en)^{n/d}\exp(\sqrt n)\notag\\
 &\le (2en)^{n/2}\exp(\sqrt n).\label{cgr:ld:eq:symbound1}
\end{align}
There are two reflection types. A reflection through gaps has no
fixed endpoints, and its fixed-matching count is exactly the
$d=2$, $k=n$ instance of \eqref{cgr:ld:eq:rotation-egf}. A reflection
through opposite endpoints has two fixed endpoints. They must be
paired to each other: a fixed endpoint paired to a nonfixed one
would acquire two partners under reflection. On the remaining
endpoints the count is the $d=2$, $k=n-1$ instance. Both satisfy
\eqref{cgr:ld:eq:symbound1}. A union bound over fewer than $4n$
nonidentity symmetries gives
\begin{equation}\label{cgr:ld:eq:symbound}
 S_n\le 4n(2en)^{n/2}\exp(\sqrt n).
\end{equation}
Stirling's formula gives $\log m_n=n\log(2n)-n+O(1)$, so the
logarithm of the right side of \eqref{cgr:ld:eq:symbound} divided by
$m_n$ is $-\tfrac12 n\log n+O(n)$. This is smaller than
$-A\log n$ for each fixed $A$ and large enough $n$.
\end{proof}

This argument explicitly covers both reflection types. It also
bounds every symmetric matching, a stronger form of exception
control than bounding only graphs with fewer than $4n$
representatives.

\subsection{Representation fibers}
For an unlabeled circle graph $G$ of order $n$, define
\[
 R_n(G)=\#\{D\in\mathcal M_n:G_D\cong G\}.
\]
This counts all indexed matching representatives, including all
inequivalent dihedral orbits. Grouping the sum by its graph gives
the exact identities
\begin{equation}\label{cgr:ld:eq:fiber}
 g_n=\sum_{D\in\mathcal M_n}\frac1{R_n(G_D)},\qquad
 c_n=\sum_{D\in\mathcal M_n}
       \frac{\one_{G_D\text{ connected}}}{R_n(G_D)}.
\end{equation}
A geometrically asymmetric matching has an orbit of size $4n$;
therefore $R_n(G_D)\ge4n$ for such a matching. On symmetric
matchings we may use the crude bound $1/R_n(G_D)\le1$.
Consequently
\begin{equation}\label{cgr:ld:eq:coarse-g}
 g_n\le h_n+S_n=O(m_n/n).
\end{equation}
Fix once and for all a constant $C$ such that
\begin{equation}\label{cgr:ld:eq:uniform-g}
 g_j\le C m_j/j\quad(j\ge1),
\end{equation}
enlarging $C$ to handle the finite initial segment. No explicit
value of this constant or its initial asymptotic onset is required.

\subsection{Deleting a bounded set of chords}
\begin{lemma}\label{cgr:ld:lem:delete-sym}
For fixed $K$ and $A>0$, the probability that some deletion of at
most $K$ chords from $D_n$ leaves a geometrically symmetric
matching is $o(n^{-A})$.
\end{lemma}
\begin{proof}
For a fixed deletion size $r\le K$, choose its $2r$ endpoints,
match them in $m_r$ ways, and choose a symmetric matching on the
remaining endpoints, with their cyclic order compressed. There
are at most
\begin{equation}\label{cgr:ld:eq:delete-bound}
 \binom{2n}{2r}m_rS_{n-r}
\end{equation}
such choices. Multiple counting only increases this bound. For
fixed $r$, the ratio $m_{n-r}/m_n$ is $O(n^{-r})$, while
$\binom{2n}{2r}=O(n^{2r})$. Lemma \ref{cgr:ld:lem:symmetry}, with a
sufficiently large exponent, makes \eqref{cgr:ld:eq:delete-bound} divided
by $m_n$ smaller than any prescribed power of $n^{-1}$. Sum over
$0\le r\le K$.
\end{proof}
This is uniform over the choice of deleted chords. Thus the
chords can later be chosen from the graph's leaf set without any
unjustified conditioning argument.
\begin{remark}[\wtag\ Credit for the method of Lemma~\ref{cgr:ld:lem:symmetry}, 5~October 2026]\label{cgr:ld:rem:symmetry-credit}
The method of Lemma~\ref{cgr:ld:lem:symmetry}, the exponential generating function \eqref{cgr:ld:eq:rotation-egf} of the matchings fixed by a rotation of order $d$ and a Cauchy bound on the circle of radius $\sqrt{k/d}$, is that of \cite[Proposition 5.14]{BBFGP24}: the proportion of matchings of size $n$ whose circle graph has fewer than $4n$ representatives is $o(n^{-n/3})$, proved by bounding the matchings fixed by a nonidentity symmetry. Neither Part cites that proposition. Its proof says that a reflection and a rotation of order two give the same enumeration sequence. That is exact for the $n$ reflections whose axis passes through gaps between endpoints. A reflection through two opposite endpoints fixes $t_{n-1}$ matchings, by the argument of Lemma~\ref{cgr:ld:lem:symmetry}, where $t_k=k![z^k]\exp(z+z^2)$ is the case $d=2$ of \eqref{cgr:ld:eq:rotation-egf}; a rotation of order two fixes $t_n$. The source's bound is unaffected, because $t_{n-1}<t_n$ for $n\ge2$.

\emph{Proof of $t_{n-1}<t_n$.} Put $f(z)=\exp(z+z^2)=\sum_ka_kz^k$, so $t_k=k!\,a_k$. From $f'=(1+2z)f$, $(k+1)a_{k+1}=a_k+2a_{k-1}$ for $k\ge1$, that is $t_{k+1}=t_k+2k\,t_{k-1}$. With $t_0=t_1=1$ every $t_k$ is positive, so $t_{k+1}>t_k$ for $k\ge1$. \qed

Lemma~\ref{cgr:ld:lem:symmetry} treats the two reflection types separately. Its conclusion $S_n=o(m_nn^{-A})$ is weaker than the source's $o(n^{-n/3})$ proportion, but \eqref{cgr:ld:eq:symbound} is explicit, and it suffices for both Parts.
\end{remark}

\section{Leaves of a typical indexed matching}\label{cgr:ld:sec:leaves}
A chord counted by $Y_n$ is called \emph{short}. Its shorter arc
contains one endpoint, so it crosses exactly the chord incident to
that endpoint. It is therefore a graph leaf. Call its unique
neighbor its parent.

\begin{lemma}\label{cgr:ld:lem:long-leaf}
The expected number of graph leaves which are not short is
$O(n^{-1})$.
\end{lemma}
\begin{proof}
A graph leaf crosses exactly one other chord. The smaller open arc
between its endpoints must contain an odd number $2r+1$ of
endpoints: one is matched outside and the other $2r$ are matched
internally. For a nonshort leaf $r\ge1$. After its two endpoints
are fixed, choose the inside and outside endpoints of the unique
crossing chord. All remaining matchings lie independently inside
or outside the arc, so the exact number of completions is
\begin{align}
 &(2r+1)(2n-2r-3)m_rm_{n-r-2}\notag\\
 &\hspace{35mm}=m_{r+1}m_{n-r-1}.\label{cgr:ld:eq:leaf-completions}
\end{align}
For any specified arc length there are at most $2n$ possible
leaf chords. Diameter choices, when present, are merely
counted twice. Put $s=r+1$ and
\begin{equation}\label{cgr:ld:eq:q-def}
 q_s=\frac{m_sm_{n-s}}{m_n}.
\end{equation}
The expectation in question is at most
\begin{equation}\label{cgr:ld:eq:long-expectation}
 2n\sum_{s=2}^{\floor{n/2}}q_s.
\end{equation}
For $0\le s\le(n-1)/2$,
\begin{equation}\label{cgr:ld:eq:q-ratio}
 \frac{q_{s+1}}{q_s}=\frac{2s+1}{2n-2s-1}\le1.
\end{equation}
Thus, for all sufficiently large $n$,
\begin{align}
 \sum_{s=2}^{\floor{n/2}}q_s
 &\le q_2+nq_3=O(n^{-2}),\label{cgr:ld:eq:q-tail}\\
 q_2&=\frac3{(2n-1)(2n-3)},\notag\\
 q_3&=\frac{15}{(2n-1)(2n-3)(2n-5)}.\notag
\end{align}
Substitution into \eqref{cgr:ld:eq:long-expectation} proves the claim.
\end{proof}

\begin{lemma}\label{cgr:ld:lem:parents}
The probability that two short graph leaves have the same parent
is $O(n^{-1})$.
\end{lemma}
\begin{proof}
Choose the two short-arc positions in $O(n^2)$ ways. When their
configuration is feasible, the two leaf chords and their common
parent are three specified distinct matching edges on six
endpoints: the parent joins the unique endpoint enclosed by each
short leaf. A uniform matching contains these three edges with
probability $m_{n-3}/m_n=O(n^{-3})$. Impossible or overlapping
configurations contribute zero. Summing bounds the expected
number of such pairs, hence also the probability, by $O(n^{-1})$.
\end{proof}

In a connected graph of order greater than two, a leaf's parent
cannot itself be a leaf. Lemmas \ref{cgr:ld:lem:long-leaf} and
\ref{cgr:ld:lem:parents} imply that, outside an unconditional event of
probability $O(n^{-1})$, every graph leaf is short and all parents
are distinct. The estimates remain adequate after we restrict a
weighted sum to connected graphs: no conditional probability
claim is needed.

\section{An upper bound from independent leaf moves}\label{cgr:ld:sec:upper}
Fix $K\ge0$. Call a connected matching good if it has at most
$K$ graph leaves, all of them are short, their parents are distinct,
and deleting all leaves leaves a geometrically asymmetric matching
$D_0$. By the previous lemmas, the bad connected matchings are
contained in $\{Y_n>K\}$ and an event of probability $o(1)$.
The geometric-symmetry contribution will be estimated separately.

\begin{lemma}[Leaf moves]\label{cgr:ld:lem:leaf-moves}
If a good matching has $k$ leaves, then its graph $G$ satisfies
\begin{equation}\label{cgr:ld:eq:fiber-lower}
 R_n(G)\ge4n\,2^k.
\end{equation}
\end{lemma}
\begin{proof}
View $D_0$ as a cyclic double-occurrence word, one symbol for each
retained chord. At either occurrence of each parent symbol $p$,
replace that occurrence by the block $\ell,p,\ell$, where
$\ell$ is its deleted leaf. This new chord crosses just $p$.
There are two choices per parent. Distinct parents have distinct
symbols, so their endpoint occurrences are distinct and the local
insertions commute, even when occurrences are consecutive.
Every one of the $2^k$ choices represents $G$.

Suppose two expanded diagrams are dihedrally equivalent. The
induced graph isomorphism preserves the full set of degree-one
vertices, which is exactly the inserted leaf set. Delete that set
on both sides. The remaining equivalence is a dihedral symmetry
of $D_0$ with its cyclic order compressed. Its geometric
asymmetry forces the identity on retained endpoints. The parent
endpoints are therefore individually fixed. Because different
leaves have different parents, each leaf must occupy the same
chosen parent endpoint in both expansions. The choices agree.
Hence the $2^k$ expanded diagrams lie in distinct dihedral orbits.

The same argument applied to a symmetry of a single expanded
diagram gives the identity on all retained endpoints. For $n-K$
large enough there are at least three retained endpoints; an
endpoint rotation or reflection fixing these individually is the
identity. Therefore every expanded diagram is itself geometrically
asymmetric and has exactly $4n$ indexed representatives in its
orbit. Their disjoint orbits prove \eqref{cgr:ld:eq:fiber-lower}.
\end{proof}

The lemma only gives a lower bound on the fiber size. It does not
need to classify additional representations; extra representations
would make the reciprocal weight still smaller.

For a connected matching with $n>1$, necessarily $X_n=0$, since
a chord joining consecutive endpoints is an isolated graph vertex.
On good matchings $k=Y_n$. Using \eqref{cgr:ld:eq:fiber},
\eqref{cgr:ld:eq:fiber-lower}, the weight $1/(4n)$ on all asymmetric
exceptions and the weight at most one on symmetric exceptions,
we obtain
\begin{equation}\label{cgr:ld:eq:upper-before-limit}
 \frac{c_n}{h_n}\le
 \EE\bigl[\one_{X_n=0}\one_{Y_n\le K}2^{-Y_n}\bigr]
 +\PP(Y_n>K)+o(1)+\frac{4nS_n}{m_n}.
\end{equation}
The last term tends to zero by Lemma \ref{cgr:ld:lem:symmetry}.
Convergence in distribution of integer-valued variables gives
convergence of each fixed finite collection of point
probabilities, and the Poisson limit is tight. Thus
\eqref{cgr:ld:eq:poi-source} implies, for fixed $K$,
\begin{align}
 \limsup_{n\to\infty}\frac{c_n}{h_n}
 &\le e^{-1}\sum_{j=0}^{K}2^{-j}\frac{e^{-1}}{j!}
       +\PP(\Poi(1)>K).\label{cgr:ld:eq:upper-truncation}
\end{align}
Letting $K\to\infty$ gives
\begin{equation}\label{cgr:ld:eq:upper-final}
 \limsup_{n\to\infty}\frac{c_n}{h_n}\le e^{-3/2}.
\end{equation}

\section{A lower bound from recoverable prime decorations}\label{cgr:ld:sec:lower}
\subsection{Asymmetric prime cores}
Let $p_N^{\mathrm a}$ count asymmetric split-prime circle graphs
of order $N$. By \eqref{cgr:ld:eq:prime-prob}, the number of prime
indexed matchings is $(e^{-3}+o(1))m_N$. Removing all geometrically
symmetric matchings changes this by at most $S_N=o(m_N)$.
Prime uniqueness and \eqref{cgr:ld:eq:prime-asym} identify every remaining
graph with one full orbit of size $4N$. Therefore
\begin{equation}\label{cgr:ld:eq:prime-count}
 p_N^{\mathrm a}\sim e^{-3}\frac{m_N}{4N}.
\end{equation}

A split-prime graph of order at least five is connected, has
minimum degree at least two, and has no twin pair. Here vertices
$u,v$ are twins if their adjacencies to every vertex outside
$\{u,v\}$ agree; they are true twins if adjacent and false twins
if nonadjacent. A twin pair is a split side of size two. A leaf
and its parent also form a split side of size two. Disconnected
graphs admit an empty crossing split, or an equivalent split
using an isolated vertex and another vertex when a component has
order one. These observations establish the stated properties.

\subsection{Three single-vertex decorations}
Start from an asymmetric split-prime core $H$ of order $N\ge5$.
Choose $k$ distinct core vertices. At each selected vertex perform
exactly one of these operations, each adding one vertex:
\begin{enumerate}
\item[$L$.] Add one pendant leaf adjacent to that vertex.
\item[$T$.] Replace the vertex by two adjacent twins, each with
its former neighbors.
\item[$F$.] Replace the vertex by two nonadjacent twins, each
with its former neighbors.
\end{enumerate}
No core vertex receives more than one decoration. The resulting
graph is connected and has $n=N+k$ vertices. All operations
preserve the circle-graph class. For $L$, insert a chord around one
endpoint as in Lemma \ref{cgr:ld:lem:leaf-moves}. For $T$ or $F$,
duplicate the two endpoints of the core chord in very small
adjacent neighborhoods. Alternating the two chord symbols makes
them cross; nesting them makes them noncrossing. Both new chords
meet exactly the former external neighbors. These local
substitutions can be made simultaneously at distinct core chords.

\begin{lemma}[Canonical recovery]\label{cgr:ld:lem:recovery}
The decorated graph determines its core isomorphism class, every
selected core vertex, and its $L$, $T$ or $F$ type. Two decoration
patterns on an asymmetric core are isomorphic only if they are
identical.
\end{lemma}
\begin{proof}
Every old or cloned vertex retains at least the core's minimum
degree two. The added $L$ vertices are consequently exactly the
degree-one vertices of the result. Delete precisely those leaves
and save their incidences with their parents.

What remains is a blowup of the twin-free graph $H$ with each
replacement block of size one or two. Distinct core vertices
$u,v$ have a third vertex $w$ whose adjacency distinguishes them,
since otherwise $u,v$ would be twins. Every replacement block
for $w$ is nonempty. Any vertex in that block continues to
distinguish a vertex in the $u$ block from a vertex in the $v$
block. Thus twins in the blowup cannot lie in different blocks.
The only nonsingleton maximal twin classes are exactly the
intended size-two $T$ and $F$ blocks.

Contract each such class to one vertex. Its internal edge or
nonedge records whether its type was $T$ or $F$. The quotient is
$H$, and the saved pendant incidences identify its $L$ vertices.
All steps are canonical under graph isomorphisms. An isomorphism
between two outputs therefore induces an isomorphism between
cores preserving all decorations. On the same asymmetric core
that isomorphism must fix every vertex. This proves the assertion.
\end{proof}

Twin swaps in the output graph are allowed, and indeed occur.
There is no extra division by two for each twin pair: the count
is of decoration patterns on the recovered asymmetric core, not
of independently labeled clones.

\subsection{Normalization and the lower sandwich}
For fixed $k$ the number of constructed graphs of order $n$ is
exactly
\begin{equation}\label{cgr:ld:eq:decorate-count}
 3^k\binom{n-k}{k}p_{n-k}^{\mathrm a}
\end{equation}
when $n-k\ge\max(5,k)$. The families for different $k$ are
disjoint, since recovery determines the core order. For every
fixed $K$ and all sufficiently large $n$,
\begin{equation}\label{cgr:ld:eq:lower-sum}
 c_n\ge\sum_{k=0}^{K}3^k\binom{n-k}{k}p_{n-k}^{\mathrm a}.
\end{equation}
For fixed $k$,
\begin{equation}\label{cgr:ld:eq:ratio-m}
 \frac{m_{n-k}}{m_n}
 =\prod_{j=0}^{k-1}\frac1{2n-2j-1}
 \sim(2n)^{-k},\qquad
 \binom{n-k}{k}\sim\frac{n^k}{k!}.
\end{equation}
Using \eqref{cgr:ld:eq:prime-count}, the exact normalization is
\begin{align}
 \frac{3^k\binom{n-k}{k}p_{n-k}^{\mathrm a}}{h_n}
 &\sim e^{-3}3^k\binom{n-k}{k}
          \frac n{n-k}\frac{m_{n-k}}{m_n}\notag\\
 &\longrightarrow e^{-3}\frac{(3/2)^k}{k!}.\label{cgr:ld:eq:lower-term}
\end{align}
First take the liminf in \eqref{cgr:ld:eq:lower-sum} with $K$ fixed,
and then increase $K$. This gives
\begin{equation}\label{cgr:ld:eq:lower-final}
 \liminf_{n\to\infty}\frac{c_n}{h_n}
 \ge e^{-3}\sum_{k=0}^{\infty}\frac{(3/2)^k}{k!}
 =e^{-3/2}.
\end{equation}
Together with \eqref{cgr:ld:eq:upper-final}, this proves the connected
equivalent. The argument never assumes a uniform estimate when
$k$ grows with $n$ and never interchanges an unbounded sum with
an unproved asymptotic expansion.

\section{Disconnected graphs and the first connectivity correction}\label{cgr:ld:sec:disconnected}
Circle graphs are closed under disjoint union, by placing their
chord diagrams in disjoint small arcs. Every disconnected graph
has a component of order $k\le n/2$. Choosing one such component
and its complement may overcount, so
\begin{equation}\label{cgr:ld:eq:component-bound}
 0\le g_n-c_n\le
 \sum_{k=1}^{\floor{n/2}}g_kg_{n-k}.
\end{equation}
No binomial coefficient occurs: these are unlabeled graph
isomorphism classes, and disjoint union is a well-defined map
from pairs of such classes.

Use \eqref{cgr:ld:eq:uniform-g} and $n-k\ge n/2$ to get
\begin{align}
 g_n-c_n
 &\le\frac{2C^2m_n}{n}
       \sum_{k=1}^{\floor{n/2}}\frac{q_k}{k}
   =O(m_n/n^2),\label{cgr:ld:eq:disconnected-small}
\end{align}
where $q_k$ is \eqref{cgr:ld:eq:q-def}, $q_1=1/(2n-1)$ and
\eqref{cgr:ld:eq:q-tail} bounds the remaining sum. Since
$c_n\sim e^{-3/2}m_n/(4n)$, this proves
$g_n/c_n=1+O(1/n)$ and hence the all-graph equivalent in
Theorem \ref{cgr:ld:thm:main}. Stirling's formula proves its elementary
factorial form.

For a sharper comparison, graphs with at least one isolated vertex
are in bijection with all graphs on $n-1$ vertices, by removing
or adding one indistinguishable isolated vertex. The image is
unique even when several isolates are present. Disconnected
graphs with no isolated vertex have a component of order between
$2$ and $n/2$, and their count is at most
\begin{equation}\label{cgr:ld:eq:no-isolates}
 \sum_{k=2}^{\floor{n/2}}g_kg_{n-k}=O(m_n/n^3).
\end{equation}
Thus $g_n-c_n=g_{n-1}+O(m_n/n^3)$. Also
\eqref{cgr:ld:eq:disconnected-small} at $n-1$ gives
$g_{n-1}-c_{n-1}=O(m_n/n^3)$. Finally,
\begin{equation}\label{cgr:ld:eq:ratio-prev}
 \frac{g_{n-1}}{g_n}
 \sim\frac{h_{n-1}}{h_n}
 =\frac{n}{(n-1)(2n-1)}\sim\frac1{2n}.
\end{equation}
Equations \eqref{cgr:ld:eq:gap} and \eqref{cgr:ld:eq:connected} follow.

\section{Saturation and the typical recovered core}\label{cgr:ld:sec:saturation}
The upper and lower bounds agree exactly in the limit. This
forces the recoverable decorated families to carry almost all
unlabeled graphs; it does not require a structural classification
of every exceptional graph.

For nonnegative integers $\ell,t,f$, put $k=\ell+t+f$ and let
$\mathcal D_{n;\ell,t,f}$ consist of constructions with exactly
$\ell$ leaves, $t$ true-twin pairs and $f$ false-twin pairs,
on distinct vertices of an asymmetric prime core of order
$N=n-k$. Canonical recovery makes all these families disjoint,
and for $N\ge\max(5,k)$ their exact size is
\begin{equation}\label{cgr:ld:eq:triple-count}
 |\mathcal D_{n;\ell,t,f}|
 =p_N^{\mathrm a}\frac{N!}{\ell!t!f!(N-k)!}.
\end{equation}
Consequently, for every fixed triple,
\begin{equation}\label{cgr:ld:eq:triple-h}
 \frac{|\mathcal D_{n;\ell,t,f}|}{h_n}
 \longrightarrow e^{-3}\frac{2^{-k}}{\ell!t!f!}.
\end{equation}
Divide by either connected or all-graph counts, using
Theorem \ref{cgr:ld:thm:main}, to get the limiting probability
\begin{equation}\label{cgr:ld:eq:triple-prob}
 e^{-3/2}\frac{(1/2)^\ell}{\ell!}
              \frac{(1/2)^t}{t!}\frac{(1/2)^f}{f!}.
\end{equation}
The sum of \eqref{cgr:ld:eq:triple-prob} over all triples is one.
Given $\eps>0$, choose a finite set of triples whose limiting
mass exceeds $1-\eps$. Convergence on that finite set proves
that the probability of the union of \emph{all} decorated
families has liminf at least $1-\eps$. Let $\eps$ decrease to
zero. This proves saturation without exchanging an infinite sum
and an asymptotic estimate.

\begin{theorem}[Typical core and decoration laws]\label{cgr:ld:thm:core}
For a uniform unlabeled circle graph of order $n$, either
conditioned to be connected or not, with probability tending to
one there is a canonically recovered asymmetric split-prime core
of the form described in Lemma \ref{cgr:ld:lem:recovery}. Its decoration
counts satisfy
\begin{equation}\label{cgr:ld:eq:three-poi}
 (L_n,T_n,F_n)\xrightarrow{d}
 \bigl(\Poi(1/2),\Poi(1/2),\Poi(1/2)\bigr),
\end{equation}
with independent coordinates. Its deficit $n-|H|$ converges in
distribution to $\Poi(3/2)$.
\end{theorem}
\begin{proof}
On the saturated family the counts and the core are uniquely
recovered, and \eqref{cgr:ld:eq:triple-prob} gives their joint point
probabilities. The finite-set argument also proves tightness,
which yields the joint convergence. Outside this family define
the counts arbitrarily; its probability tends to zero and cannot
change the limit. Summing the three independent Poisson limits
gives the deficit law.
\end{proof}
The core claim concerns the explicit leaf-deletion and
maximal-twin-contraction procedure on a probability-one limiting
family. It is not an assertion that arbitrary iterative deletion
and contraction orders give the same core for every graph.

\subsection{Automorphisms and a labeled consequence}\label{cgr:ld:sub:aut}
\begin{proposition}\label{cgr:ld:prop:aut}
For a graph in $\mathcal D_{n;\ell,t,f}$,
\begin{equation}\label{cgr:ld:eq:aut-group}
 \Aut(G)\cong (C_2)^{t+f}.
\end{equation}
Thus a uniform unlabeled graph in either model is asymmetric
with probability tending to $e^{-1}$, and its automorphism-group
order converges in distribution to $2^P$, where $P\sim\Poi(1)$.
\end{proposition}
\begin{proof}
Every automorphism preserves the leaves, the recovered twin
classes, and the recovered core with its decorations. The core
is asymmetric, so each core vertex is fixed. A leaf then is fixed
by its distinct fixed parent. Each twin pair can be swapped
independently; these swaps are automorphisms and are the only
remaining freedoms. This proves \eqref{cgr:ld:eq:aut-group}. Saturation
and Theorem \ref{cgr:ld:thm:core} give $T_n+F_n\to\Poi(1)$ and the
asserted consequences. The arbitrary exceptional graphs have
vanishing probability.
\end{proof}
In particular, geometric asymmetry of typical representations
must not be used to claim asymmetry of typical graphs. The twin
swaps above generally do not act as rotations or reflections of
a given nonprime chord representation.

Other direct consequences are
\begin{equation}\label{cgr:ld:eq:other-prob}
 \PP(G\text{ leafless})\longrightarrow e^{-1/2},\qquad
 \PP(G\text{ split-prime})\longrightarrow e^{-3/2}.
\end{equation}
The first follows from $L_n=0$. For the second, every nonempty
decoration creates a nontrivial split: use a twin pair, or a
leaf-parent pair, as a side of size two. The zero-decoration
family is prime. Saturation accounts for all remaining cases.

Let $g_n^{\mathrm{lab}}$ and $c_n^{\mathrm{lab}}$ count graphs
on a fixed labeled $n$-element vertex set. The number of labelings
of an unlabeled graph is $n!/|\Aut(G)|$. Hence
\begin{equation}\label{cgr:ld:eq:label-expectation}
 \frac{g_n^{\mathrm{lab}}}{n!g_n}
 =\EE_{\mathrm{unlabeled}}\left[\frac1{|\Aut(G)|}\right]
 \longrightarrow\EE[2^{-P}]=e^{-1/2}.
\end{equation}
The expectation passage is legitimate because the reciprocal
automorphism size is bounded by one: first truncate to finitely
many decoration triples, then use saturation to bound the
remainder. No uniform integrability assertion about $|\Aut(G)|$
is being made. The same argument holds in the connected model.
Therefore
\begin{equation}\label{cgr:ld:eq:labeled}
 c_n^{\mathrm{lab}}\sim g_n^{\mathrm{lab}}
 \sim e^{-2}\frac{n!m_n}{4n}
 =e^{-2}\frac{(2n)!}{2^{n+2}n}.
\end{equation}
This is a mathematical consequence of the proved unlabeled law.
A separate literature and OEIS identification audit for labeled
counts was not completed, so no separate priority claim is made.
\begin{writenote}[Refined in Part~II, 5~October 2026]
Part~II's Section~\ref{cgr:fc:sec:weighted} gives the first corrections of these laws: $\PP(\Aut(G)=1)=e^{-1}(1-\frac1{2n}+O(n^{-2}))$ \eqref{cgr:fc:eq:asymmetric}, $\EE|\Aut(G)|^{-1}=e^{-1/2}(1-\frac1{6n}+O(n^{-2}))$, and the labeled counts \eqref{cgr:fc:eq:labeledc}--\eqref{cgr:fc:eq:labeledg} with relative corrections $-\frac{47}{12n}$ (connected) and $-\frac{41}{12n}$ (all). Part~II's $C_n^{\mathrm{lab}}$, $G_n^{\mathrm{lab}}$ are this Part's $c_n^{\mathrm{lab}}$, $g_n^{\mathrm{lab}}$, and $n!e^{-2}h_n=e^{-2}(2n)!/(2^{n+2}n)$ is the leading term of \eqref{cgr:ld:eq:labeled}. The leafless and split-prime laws \eqref{cgr:ld:eq:other-prob} are this Part's alone. The labeled audit called for here was not made at the write either (Section~\ref{cgr:sec:further}).
\end{writenote}

\section{Quantitative error bounds}\label{cgr:ld:sec:quantitative}
The preceding proof establishes the leading equivalent using the
published qualitative Poisson limit. We now prove a quantitative
local approximation directly and propagate it through the same
sandwich. This gives the following strengthening.

\begin{theorem}\label{cgr:ld:thm:quantitative}
For $a_n=c_n$ or $a_n=g_n$,
\begin{equation}\label{cgr:ld:eq:quant-main}
 a_n=e^{-3/2}h_n\bigl(1+O(n^{-1})\bigr),\qquad
 \frac{c_n}{g_n}=1-\frac1{2n}+O(n^{-2}).
\end{equation}
The exceptional probability outside the canonically decorated-core
family is $O(n^{-1})$. Its three defect counts are within
$O(n^{-1})$ total variation of independent $\Poi(1/2)$ variables,
and its core deficit is within $O(n^{-1})$ total variation of
$\Poi(3/2)$. The asymmetry, leafless and split-prime probabilities
in Section \ref{cgr:ld:sec:saturation} have additive error $O(n^{-1})$;
the labeled equivalent \eqref{cgr:ld:eq:labeled} has relative error
$O(n^{-1})$.
\end{theorem}
The constants and sufficiently large onsets for graph counts are
not made explicit. Only the local-pattern bound below has a
stated numerical constant. No individual first correction
coefficient is identified.
\begin{writenote}[Superseded in part, 5~October 2026]
For the counts, Part~II's Theorem~\ref{cgr:fc:thm:main} identifies the first relative coefficients ($-\frac{15}4$ and $-\frac{13}4$ in the $h_n$ normalization) and so implies the first estimate of \eqref{cgr:ld:eq:quant-main}; its Theorem~\ref{cgr:fc:thm:transfer-summary} sharpens the second. The total-variation statements, the leafless and split-prime probabilities with error $O(n^{-1})$, and the explicit local bound \eqref{cgr:ld:eq:explicit-TV} are this Part's alone.
\end{writenote}

\subsection*{Atomic patterns and an exact forcing coupling}
Put $M=2n\ge12$, with endpoints modulo $M$. A prescribed
matching pattern of $r$ edges occurs with probability
\begin{equation}\label{cgr:ld:eq:pattern-prob}
 \rho_r=\prod_{j=0}^{r-1}(M-2j-1)^{-1}.
\end{equation}
Take as atomic $X$ patterns the $M$ edges $\{i,i+1\}$, as
$Y$ patterns the $M$ edges $\{i,i+2\}$, and as $Z$ patterns
the two possible matchings between each unordered pair of
disjoint cycle edges. There are $M(M-3)/2$ pairs of disjoint
cycle edges, hence $M(M-3)$ two-edge $Z$ patterns.

These two-edge patterns are distinct for $M\ge6$. Otherwise the
same four endpoints would admit two distinct perfect matchings
by cycle edges; their symmetric difference would be a four-cycle
inside $C_M$, which is impossible. Summing pattern indicators
therefore gives exactly the source variables $X_n,Y_n,Z_n$.
Their means are all equal to
\begin{equation}\label{cgr:ld:eq:local-means}
 \lambda_X=\lambda_Y=M\rho_1=\frac M{M-1},\qquad
 \lambda_Z=M(M-3)\rho_2=\frac M{M-1}=:\lambda.
\end{equation}
Let $\mathcal I$ be the finite pattern set, let $I_\alpha$ be
the occurrence indicator of pattern $\alpha$, and write
$p_\alpha=\EE I_\alpha$ and $S_\alpha$ for its endpoint support.

To force an edge $\{u,v\}$ into a matching, do nothing if it is
already there. Otherwise replace $\{u,x\},\{v,y\}$ by
$\{u,v\},\{x,y\}$. To force a two-edge pattern, apply this
operation in a fixed order, excluding already forced endpoints
on the next step. Denote the result by $D^\alpha$.

The result has precisely the uniform law conditional on
$I_\alpha=1$. On $L$ currently unforced endpoints, each
matching containing a prescribed edge has $L-1$ preimages:
itself and two preimages for each of its $L/2-1$ other edges.
One force thus pushes the uniform matching law to the uniform
conditional law. Applying this fact successively proves the
assertion for two edges.

Every original matching edge whose endpoints lie outside
$S_\alpha$ survives every forcing step. Only edges incident
with a forced endpoint are removed. Therefore
\begin{equation}\label{cgr:ld:eq:forcing-monotone}
 S_\alpha\cap S_\beta=\varnothing
 \quad\Longrightarrow\quad I_\beta(D^\alpha)\ge I_\beta(D).
\end{equation}
This is a statement about disjoint supports only.

\subsection*{A self-contained Poisson approximation bound}
We use total variation in the convention
$d_{\rm TV}(P,Q)=\sup_B|P(B)-Q(B)|$.

\begin{lemma}[Reduced conditional coupling bound]\label{cgr:ld:lem:stein}
Let $\Xi=(I_\alpha)_{\alpha\in\mathcal I}$ be any finite family
of indicators with means $p_\alpha$. Suppose $\Xi^\alpha$ has
the law of $\Xi-\delta_\alpha$ conditional on $I_\alpha=1$.
If $\Pi$ has independent Poisson coordinates of means
$p_\alpha$, then
\begin{equation}\label{cgr:ld:eq:stein-bound}
 d_{\rm TV}(\mathcal L(\Xi),\mathcal L(\Pi))
 \le\sum_\alpha p_\alpha\EE\|\Xi-\Xi^\alpha\|_1.
\end{equation}
\end{lemma}
\begin{proof}
For a function $f$ on nonnegative integer vectors write
$\Delta_\alpha f(\xi)=f(\xi+\delta_\alpha)-f(\xi)$.
The independent immigration--death process has generator
\[
 \mathcal Af(\xi)=\sum_\alpha p_\alpha\Delta_\alpha f(\xi)
 -\sum_\alpha\xi_\alpha\Delta_\alpha f(\xi-\delta_\alpha).
\]
Each initial particle dies at rate one; particles immigrate in
coordinate $\alpha$ at rate $p_\alpha$. Its equilibrium law is
$\Pi$. For an event indicator $J$, the semigroup solution
\[
 f_J(\xi)=-\int_0^\infty
 \bigl(\EE J(Z_\xi(t))-\EE J(\Pi)\bigr)\,dt
\]
satisfies $\mathcal Af_J=J-\EE J(\Pi)$. The integral is finite:
couple equilibrium and nonequilibrium initial states with the
same immigrants, so that their finite expected initial
population difference survives with probability decaying as
$e^{-t}$.

Couple processes with zero, one and two additional initial
particles, including two in the same coordinate if needed. A
mixed second difference of a function valued in $[0,1]$ has
absolute value at most two; it vanishes unless both additional
particles survive. Every second difference of $f_J$ is thus
bounded in absolute value by
$\int_0^\infty2e^{-2t}\,dt=1$. By telescoping coordinate
changes, $\Delta_\alpha f_J$ is 1-Lipschitz in $\ell^1$.
Conditional expectation of the death term now gives
\[
 \EE\mathcal Af_J(\Xi)=\sum_\alpha p_\alpha\EE
 \bigl[\Delta_\alpha f_J(\Xi)-
       \Delta_\alpha f_J(\Xi^\alpha)\bigr].
\]
Taking absolute values, then the supremum over events $J$,
proves \eqref{cgr:ld:eq:stein-bound}.
\end{proof}

Use $\Xi^\alpha_\alpha=0$ and
$\Xi^\alpha_\beta=I_\beta(D^\alpha)$ for $\beta\ne\alpha$.
The coordinate $\alpha$ contributes $p_\alpha^2$ to
\eqref{cgr:ld:eq:stein-bound}. By \eqref{cgr:ld:eq:forcing-monotone}, a
pair on disjoint supports contributes exactly
\[
 \PP(I_\alpha=I_\beta=1)-p_\alpha p_\beta.
\]
An overlapping pair contributes at most
$p_\alpha p_\beta+\PP(I_\alpha=I_\beta=1)$. We bound the
diagonal, overlapping product, compatible overlap, and disjoint
sums separately, always using ordered pairs when summing over
$\alpha,\beta$.

The diagonal sum is
\begin{equation}\label{cgr:ld:eq:B0}
 B_0=\sum_\alpha p_\alpha^2
 =\frac{2M}{(M-1)^2}+\frac M{(M-1)^2(M-3)}.
\end{equation}
At an endpoint $v$, cyclic symmetry and the support sizes give
\[
 \mu_v=\sum_{\alpha:v\in S_\alpha}p_\alpha
 =\frac{2\lambda_X+2\lambda_Y+4\lambda_Z}{M}
 =\frac8{M-1}.
\]
Thus the product contribution from overlapping pairs is at most
\begin{equation}\label{cgr:ld:eq:B1}
 B_1\le\sum_v\mu_v^2=\frac{64M}{(M-1)^2}.
\end{equation}
Including diagonal pairs in this upper bound is harmless.

Two compatible matching patterns sharing an endpoint must share
its incident chord. Distinct singleton patterns cannot overlap
compatibly. Every fixed chord belongs to at most four $Z$
patterns: choose one cycle neighbor of each of its endpoints,
after which the second chord is determined. There are $2M$
singleton patterns, so there are at most $16M$ ordered
compatible singleton/$Z$ pairs. Each union forces two edges.
Two distinct compatible $Z$ patterns cannot share both chords,
by the no-duplication argument. Choose their unique shared chord
in at most $\binom M2$ ways and an ordered pair of distinct
$Z$ patterns containing it in at most $4\cdot3$ ways. Their
union forces three edges. Therefore
\begin{align}
 B_2&\le16M\rho_2+6M(M-1)\rho_3\notag\\
 &=\frac{16M}{(M-1)(M-3)}
       +\frac{6M}{(M-3)(M-5)}.\label{cgr:ld:eq:B2}
\end{align}
Incompatible pairs have zero joint occurrence probability.

For disjoint patterns with $r,s\in\{1,2\}$ edges, the joint
probability is $\rho_{r+s}$ and
\[
 1\le\frac{\rho_{r+s}}{\rho_r\rho_s}
 =\prod_{j=0}^{s-1}\frac{M-2j-1}{M-2r-2j-1}
 \le\frac{(M-1)(M-3)}{(M-5)(M-7)}.
\]
The largest ratio occurs at $r=s=2$. With
$\delta_M=8(M-4)/((M-5)(M-7))$, the total disjoint
contribution is consequently bounded by
\begin{equation}\label{cgr:ld:eq:B3}
 B_3\le\delta_M\left(\sum_\alpha p_\alpha\right)^2
 =\delta_M\left(\frac{3M}{M-1}\right)^2.
\end{equation}
All four bounds are $O(M^{-1})$. Summing Poisson pattern
coordinates by their type gives independent Poisson variables
of mean $\lambda$. Coupling each to $\Poi(1)$ costs at most
$\lambda-1$, by adding an independent Poisson increment.
Total variation decreases under the summing map, so
\begin{equation}\label{cgr:ld:eq:local-TV}
 d_{\rm TV}\bigl(\mathcal L(X_n,Y_n,Z_n),\Poi(1)^{\otimes3}\bigr)
 \le B_0+B_1+B_2+B_3+\frac3{M-1}=O(M^{-1}).
\end{equation}
For a numerical bound, multiply the explicit right side by
$M$. Every summand decreases for $M\ge12$; for the only
less immediate term, use
\[
 MB_3\le72\left(\frac M{M-1}\right)^2
              \frac M{M-5}\frac{M-4}{M-7}.
\]
At $M=12$ their sum is less than $354$, hence less than $400$.
We have proved, for $n\ge6$,
\begin{equation}\label{cgr:ld:eq:explicit-TV}
 d_{\rm TV}\bigl(\mathcal L(X_n,Y_n,Z_n),\Poi(1)^{\otimes3}\bigr)
 \le\min\{1,200/n\}.
\end{equation}
The generous constant is not intended as a sharp small-$n$
approximation.
\begin{writenote}[Checked, 5~October 2026]
In exact rational arithmetic, $M(B_0+B_1+B_2+B_3+\frac3{M-1})$ with the right-hand sides of \eqref{cgr:ld:eq:B0}--\eqref{cgr:ld:eq:B3} is $353.96$ at $M=12$ (rounded), and $299.25$, $236.50$, $192.12$ at $M=14$, $20$, $40$.
\end{writenote}

\subsection*{Prime cores and a logarithmic cutoff}
The published equivalence between simultaneous vanishing of the
local counts and absence of the endpoint decomposition sizes,
together with \eqref{cgr:ld:eq:large-decomp}, implies
\[
 \left|\PP(D_n\text{ indecomposable})-
       \PP(X_n=Y_n=Z_n=0)\right|=O(n^{-1}).
\]
Using \eqref{cgr:ld:eq:local-TV} and the symmetry bound, the argument
of \eqref{cgr:ld:eq:prime-count} therefore improves to
\begin{equation}\label{cgr:ld:eq:quant-prime}
 p_n^{\mathrm a}=e^{-3}h_n\bigl(1+O(n^{-1})\bigr).
\end{equation}
Its implicit constant is independent of its integer argument,
so it is uniform for arguments in
$[n-\ceil{\log n},n]$.

Put $K=\ceil{\log n}$, with natural logarithms. Write
$(Y)_K=Y(Y-1)\cdots(Y-K+1)$ for the falling factorial.
Any compatible ordered list of $K$ distinct $Y$ indicators
forces $K$ distinct matching edges. Thus
\begin{align}
 \EE(Y_n)_K
 &\le(2n)^K\frac{m_{n-K}}{m_n}\notag\\
 &=\prod_{j=0}^{K-1}\frac{2n}{2n-2j-1}
 \le\exp\{O(K^2/n)\}.\label{cgr:ld:eq:Y-factorial}
\end{align}
Since $(Y_n)_K\ge K!$ on $Y_n\ge K$,
\begin{equation}\label{cgr:ld:eq:Y-tail-quant}
 \PP(Y_n>K)\le\frac{\exp\{O(K^2/n)\}}{K!}
   =o(n^{-a})
\end{equation}
for every fixed $a>0$. This is a direct growing-order estimate,
not an unsupported use of fixed-moment convergence.

The deletion bound \eqref{cgr:ld:eq:delete-bound} also extends to
$r\le K$. Uniformly over these $r$,
\[
 \binom{2n}{2r}m_r\le\exp\{O((\log n)^2)\},
 \qquad m_{n-r}/m_n\le1.
\]
By \eqref{cgr:ld:eq:symbound},
$S_N/m_N\le\exp\{-\tfrac12N\log N+O(N)\}$.
Apply it at $N=n-r$ and sum over $r\le K$ to obtain
\begin{equation}\label{cgr:ld:eq:delete-quant}
 \PP(\text{some deletion of at most }K\text{ chords is symmetric})
 \le\exp\{-\tfrac12n\log n+O(n+(\log n)^2)\}.
\end{equation}
This remains smaller than every fixed inverse power of $n$.
The bounds for long leaves and repeated parents in
Section \ref{cgr:ld:sec:leaves} already have order $O(n^{-1})$.

\subsection*{The quantitative sandwich}
Use the definition of good connected matchings with the new
$K=\ceil{\log n}$. Equations \eqref{cgr:ld:eq:Y-tail-quant} and
\eqref{cgr:ld:eq:delete-quant}, and the leaf regularity estimates, show
that \emph{connected nongood matchings} have unconditional
probability $O(n^{-1})$. The leaf-movement proof applies
unchanged, since $n-K\to\infty$. Consequently
\begin{align}
 \frac{c_n}{h_n}
 &\le\EE[\one_{X_n=0}2^{-Y_n}]+O(n^{-1})
           +\frac{4nS_n}{m_n}\notag\\
 &=e^{-3/2}+O(n^{-1}).\label{cgr:ld:eq:quant-upper}
\end{align}
The last equality follows from \eqref{cgr:ld:eq:local-TV}, because
the integrand lies in $[0,1]$.

For the lower bound, canonical recovery still gives
\eqref{cgr:ld:eq:lower-sum}, now for this growing $K$. The exact ratio
for $0\le k\le K$ is
\begin{align}
 \frac{3^k\binom{n-k}{k}h_{n-k}}{h_n}
 &=\frac{(3/2)^k}{k!}\frac n{n-k}
   \prod_{j=0}^{k-1}
       \frac{1-(k+j)/n}{1-(j+1/2)/n}.\label{cgr:ld:eq:uniform-product}
\end{align}
For large enough $n$ every factor is positive and its logarithm
is bounded using $|\log(1-z)|\le2|z|$ for $|z|\le1/2$.
The factors after $(3/2)^k/k!$ are therefore
$\exp\{O((k+1)^2/n)\}=1+O((k+1)^2/n)$ uniformly for
$k\le K$. Combining with \eqref{cgr:ld:eq:quant-prime} yields
\begin{equation}\label{cgr:ld:eq:uniform-decorate}
 \frac{3^k\binom{n-k}{k}p_{n-k}^{\mathrm a}}{h_n}
 =e^{-3}\frac{(3/2)^k}{k!}
       \left(1+O\left(\frac{(k+1)^2}{n}\right)\right).
\end{equation}
Importantly, sum this error against the factorial weights:
\[
 \sum_{k\ge0}\frac{(3/2)^k}{k!}(k+1)^2<\infty,\qquad
 \sum_{k>K}\frac{(3/2)^k}{k!}=o(n^{-a})
\]
for every fixed $a>0$. Hence
\begin{equation}\label{cgr:ld:eq:quant-lower}
 \frac{c_n}{h_n}\ge e^{-3/2}-O(n^{-1}).
\end{equation}
The upper and lower inequalities give the connected estimate
in \eqref{cgr:ld:eq:quant-main}. Equation \eqref{cgr:ld:eq:disconnected-small}
gives the same estimate for $g_n$. Finally,
\[
 \frac{g_{n-1}}{g_n}
 =\frac n{(n-1)(2n-1)}\bigl(1+O(n^{-1})\bigr)
 =\frac1{2n}+O(n^{-2}),
\]
and dividing \eqref{cgr:ld:eq:gap} by $g_n$ proves the strengthened
connectivity correction.

\subsection*{Quantitative saturation and automorphisms}
For $k=\ell+t+f\le K$, the exact family count
\eqref{cgr:ld:eq:triple-count}, the same product bound and the now
quantitative denominator give, uniformly over these triples,
\begin{equation}\label{cgr:ld:eq:quant-triples}
 \PP(G\in\mathcal D_{n;\ell,t,f})
 =e^{-3/2}\frac{2^{-k}}{\ell!t!f!}
       \left(1+O\left(\frac{(k+1)^2}{n}\right)\right).
\end{equation}
This holds for either uniform connected or uniform all graphs.
Sum the absolute errors over $k\le K$. The limiting deficit is
$\Poi(3/2)$ and has a finite second moment, so the sum is
$O(n^{-1})$. Its tail beyond $K$ is superpolynomially small.
Thus the probability of the complement of the truncated family,
and in particular the complement of the full decorated family,
is $O(n^{-1})$.

Give atypical graphs any defect-count convention, or a cemetery
value. Their total probability is $O(n^{-1})$; the tail and
summed errors prove the stated total-variation bounds for defect
counts and core deficit. On the family,
$\Aut(G)\cong(C_2)^{t+f}$ exactly, so bounded indicators give
\[
 \PP(G\text{ asymmetric})=e^{-1}+O(n^{-1}),\quad
 \PP(G\text{ leafless})=e^{-1/2}+O(n^{-1}),
\]
and $\PP(G\text{ split-prime})=e^{-3/2}+O(n^{-1})$.
The group-order distribution is within $O(n^{-1})$ total
variation of $2^P$, $P\sim\Poi(1)$. These assertions do not
claim convergence of unbounded automorphism moments.
The bounded reciprocal automorphism factor similarly gives
\[
 \frac{g_n^{\mathrm{lab}}}{n!g_n}=e^{-1/2}+O(n^{-1}).
\]
Multiply by the quantitative unlabeled equivalent to obtain
\begin{equation}\label{cgr:ld:eq:quant-labeled}
 g_n^{\mathrm{lab}},\ c_n^{\mathrm{lab}}
 =e^{-2}\frac{(2n)!}{2^{n+2}n}\bigl(1+O(n^{-1})\bigr),
\end{equation}
where the displayed estimate holds for each sequence separately.
This completes the proof of Theorem \ref{cgr:ld:thm:quantitative}.

\section{Smooth inversion and ceiling-safe thresholds}\label{cgr:ld:sec:inverse}
\subsection{A precise smooth model}
For either $a_n=c_n$ or $a_n=g_n$, define the positive real model
\begin{equation}\label{cgr:ld:eq:H}
 H(t)=\frac{e^{-3/2}2^t\Gamma(t+1/2)}{4\sqrt\pi\,t},
 \qquad t>0.
\end{equation}
Then $H(n)=e^{-3/2}h_n$ and $a_n/H(n)\to1$.
For large $t$, Stirling's expansion gives
\begin{align}
 \log H(t)
 &=t\bigl(\log(2t)-1\bigr)-\log t-B
       -\frac1{24t}+O(t^{-3}),\label{cgr:ld:eq:logH}\\
 B&=\frac32(1+\log2).\notag
\end{align}
The derivative is $\log(2t)-1/t+O(t^{-2})$, so $H$ is
strictly increasing on a sufficiently far real tail. Let $\phi$
be that eventual real inverse.

For $x$ large, write
\begin{equation}\label{cgr:ld:eq:u}
 L=\log x,\qquad u=\frac{L}{W_0(2L/e)},\qquad A=\log(2u),
\end{equation}
where $W_0$ is the positive principal Lambert function. The exact
identity $W_0(2L/e)e^{W_0(2L/e)}=2L/e$ implies
\begin{equation}\label{cgr:ld:eq:core-root}
 u\bigl(\log(2u)-1\bigr)=L.
\end{equation}

\begin{proposition}[Smooth inverse]\label{cgr:ld:prop:inverse}
As $x\to\infty$,
\begin{equation}\label{cgr:ld:eq:smooth-inverse}
 \phi(x)=u+1+\frac{3+\log2}{2\log(2u)}
       +O\left(\frac1{u\log u}\right).
\end{equation}
\end{proposition}
\begin{proof}
Put $F(t)=t(\log(2t)-1)$, so $F'(u)=A$ and
$F''(u)=1/u$. By \eqref{cgr:ld:eq:logH} and \eqref{cgr:ld:eq:core-root},
\[
 \log H(u)=L-\log u-B+O(u^{-1}).
\]
The derivative of $\log H$ is comparable to $A$. Evaluating
at $u$ and $u+2$ shows that the root $\phi(x)=u+d$ has
bounded displacement, eventually $0<d<2$. Uniformly for
bounded $d$, Taylor expansion gives
\begin{equation}\label{cgr:ld:eq:inverse-taylor}
 0=\log H(u+d)-L=dA-\log u-B+O(u^{-1}).
\end{equation}
Since $\log u=A-\log2$,
\[
 d=\frac{\log u+B}{A}+O((uA)^{-1})
   =1+\frac{B-\log2}{A}+O((uA)^{-1}).
\]
Finally $B-\log2=(3+\log2)/2$, proving the formula.
The $-1/(24u)$ term in \eqref{cgr:ld:eq:logH} affects only the stated
remainder, not the displayed coefficient.
\end{proof}
The shift $+1$ is forced by the factor $1/n$ in the forward
model. It is not an optional change of the OEIS offset.
\begin{writenote}[Instances of the transseries apparatus; refined in Part~II, 5~October 2026]
The equation \eqref{cgr:ld:eq:core-root} is the factorial core \texttt{p0:prop:factorial-core} of the transseries volume \cite{TSvol} with $\kappa=1$ and $d=\log2-1$, and $u$ is its solution; after taking logarithms it is the Lambert core \texttt{p0:thm:lambert-core} with $a=b=1$ (front matter, ``The OEIS entries, and the inverses''). The displacement $1+\frac{3+\log2}{2\log(2u)}$ is derived here from the Gamma model, not from the volume. Part~II's Theorem~\ref{cgr:fc:thm:inverse} carries the expansion one order further: with Part~II's $D=\log(2u)$ (this Part's $A$) and $B=\frac{3+\log2}2$ (this Part's $B-\log2$), the centre \eqref{cgr:ld:eq:s} is $u+1+B/D$, and Part~II adds $(K_a-B^2/(2D^2))/(uD)$.
\end{writenote}

\subsection{The least integer index}
Define
\begin{equation}\label{cgr:ld:eq:N}
 N_a(x)=\min\{n\ge1:a_n\ge x\}.
\end{equation}
The leading equivalent implies
\begin{equation}\label{cgr:ld:eq:successive}
 \frac{a_{n+1}}{a_n}\sim\frac{H(n+1)}{H(n)}
  =(2n+1)\frac n{n+1}\sim2n,
\end{equation}
so $a_n$ is eventually strictly increasing and unbounded.
For $x$ exceeding the maximum on a finite initial segment,
$N_a(x)$ lies in that monotone tail.

\begin{proposition}[Fixed-error ceiling bracket]\label{cgr:ld:prop:fixed-bracket}
For each fixed $\eta>0$ and all sufficiently large $x$,
\begin{equation}\label{cgr:ld:eq:eta-bracket}
 \ceil{\phi(xe^{-\eta})}\le N_a(x)
       \le\ceil{\phi(xe^{\eta})}.
\end{equation}
\end{proposition}
\begin{proof}
On a sufficiently far integer tail,
$e^{-\eta}H(n)\le a_n\le e^{\eta}H(n)$. At $n=N_a(x)$,
$a_n\ge x$ implies $H(n)\ge xe^{-\eta}$, giving the lower
bound. At $m=\ceil{\phi(xe^{\eta})}$, monotonicity of $H$
gives $H(m)\ge xe^{\eta}$ and hence $a_m\ge x$. This proves
the upper bound, once both indices are in the stated tail.
\end{proof}

A shrinking bracket can be made explicit as an \emph{existence}
statement without assuming an unproved convergence rate. Let
\begin{equation}\label{cgr:ld:eq:tail-envelope}
 r_n=\left|\log\frac{a_n}{H(n)}\right|,\qquad
 \eta(x)=\sup_{n\ge\floor{u/2}} r_n.
\end{equation}
For large enough $x$ this tail supremum is finite. Theorem
\ref{cgr:ld:thm:quantitative} gives $r_n=O(n^{-1})$ and therefore
$\eta(x)=O(u^{-1})$. It is defined from the actual sequence;
its implicit constant and onset are not supplied as computable
finite-input bounds.

To justify applying this moving envelope, first apply
Proposition \ref{cgr:ld:prop:fixed-bracket} with $\eta=1$.
Proposition \ref{cgr:ld:prop:inverse} and the derivative estimate show
$\phi(xe^{\pm1})=u+O(1)$. Thus $N_a(x)$ is eventually larger
than $u/2$, and both model endpoints formed with
$0\le\eta(x)\le1$ also lie above $u/2$. The proof of
Proposition \ref{cgr:ld:prop:fixed-bracket} now applies using the
uniform bound \eqref{cgr:ld:eq:tail-envelope}, without a circular
choice of its starting index.

By the mean value theorem applied to $\log H$ near $u$,
\begin{equation}\label{cgr:ld:eq:inverse-shift}
 \phi(xe^{\pm\eta(x)})
 =\phi(x)+O\left(\frac{\eta(x)}{\log(2u)}\right).
\end{equation}
Set
\begin{equation}\label{cgr:ld:eq:s}
 s(x)=u+1+\frac{3+\log2}{2\log(2u)}.
\end{equation}
There exists a nonnegative function $\eps(x)=O(u^{-1})$ such that,
for all sufficiently large $x$,
\begin{equation}\label{cgr:ld:eq:shrinking-bracket}
 \ceil{s(x)-\frac{\eps(x)}{\log(2u)}}
 \le N_a(x)\le
 \ceil{s(x)+\frac{\eps(x)}{\log(2u)}}.
\end{equation}
For example, sufficiently large fixed constants in
$\eps(x)=C_1\eta(x)+C_2/u$ absorb both
\eqref{cgr:ld:eq:inverse-shift} and the model remainder.

Equivalently, there are constants $C_*>0$ and $x_*>0$ such that
for $x\ge x_*$,
\begin{equation}\label{cgr:ld:eq:quant-inverse}
 \ceil{s(x)-\frac{C_*}{u\log(2u)}}
 \le N_a(x)\le
 \ceil{s(x)+\frac{C_*}{u\log(2u)}}.
\end{equation}
The width uses both the quantitative forward remainder and the
$O(1/(u\log u))$ error of the displayed smooth-model expansion;
no smaller truncation error is silently assumed.
Eventually the bracket contains at most two adjacent integers.
If its interval contains no integer, the two ceilings agree.
Near an integer it may retain two possibilities; the leading
asymptotic alone cannot resolve them uniformly. In particular,
\eqref{cgr:ld:eq:shrinking-bracket} is not a license for unconditional
nearest-integer rounding, and the report gives no effective
onset or numerical bound for $\eps(x)$.
\begin{writenote}[Refined in Part~II, 5~October 2026]
Part~II's \eqref{cgr:fc:eq:inverse-bracket} narrows the half-width $C_*/(u\log(2u))$ of \eqref{cgr:ld:eq:quant-inverse} to $C_a/(u^2\log(2u))$, with the same qualifications: noneffective constants and onset, and no unconditional rounding. These brackets follow the separation pattern of the transseries volume's \texttt{p0:thm:staircase}\,(2) but are not instances of it, since $H$ does not interpolate $a_n$ (front matter).
\end{writenote}

\section{Exact finite checks and what they establish}\label{cgr:ld:sec:checks}
\subsection{Source counts and scale}
Table \ref{cgr:ld:tab:counts} records the inspected exact initial values.
The companion reconciles these with the unlabeled component
product
\begin{equation}\label{cgr:ld:eq:Euler-product}
 \sum_{n\ge0}g_nz^n=\prod_{j\ge1}(1-z^j)^{-c_j},
\end{equation}
interpreted as a formal power series. It follows because a graph
is a multiset of connected components, uniquely up to order.
Finite truncation gives exact coefficient tests; convergence of
the displayed ordinary generating series is not assumed.

\begin{table}[htbp]
\centering
\caption{Exact inspected graph counts by order. The order-thirteen
all-graph value is from OEIS; no connected value at that order is
claimed from the inspected table.}\label{cgr:ld:tab:counts}
\begin{tabular}{rrr}
\toprule
$n$&$c_n$&$g_n$\\\midrule
1&1&1\\2&1&2\\3&2&4\\4&6&11\\5&21&34\\
6&110&154\\7&789&978\\8&8336&9497\\
9&117283&127954\\10&2026331&2165291\\
11&40302425&42609994\\12&892278076&937233306\\
13&---&22576188846\\\bottomrule
\end{tabular}
\end{table}

Inverting the logarithmic derivative of \eqref{cgr:ld:eq:Euler-product}
using the inspected all-graph value at order thirteen gives
$c_{13}=21593488017$. This is a derived value, checked by the
companion, rather than an entry in the inspected connected source
table.

For orientation, the connected ratios $c_n/h_n$ at
$n=8,9,10,11,12$ are approximately
$0.131598$, $0.122526$, $0.123797$, $0.128974$, and
$0.135436$. The all-graph ratio at $n=13$ is approximately
$0.148493$. The limiting constant is
$e^{-3/2}\approx0.223130$. These small values have not converged
rapidly enough to estimate the constant reliably. They neither
prove nor contradict the theorem.
\begin{writenote}[Reproduced, 5~October 2026]
Euler inversion of \eqref{cgr:ld:eq:Euler-product} with the terms of Table~\ref{cgr:ld:tab:counts} gives $c_{13}=21{,}593{,}488{,}017$ (also reproduced at intake). The value is not in A156808 (revision \#9), and nothing was submitted to the OEIS. With it, $c_{13}/h_{13}\approx0.142029$. Against Part~II's first corrections, $c_n/(e^{-3/2}h_n(1-\frac{15}{4n}))\approx0.8829$, $0.8946$ and $g_n/(e^{-3/2}h_n(1-\frac{13}{4n}))\approx0.8744$, $0.8873$ at $n=12$, $13$: like the ratios above, these small orders are far from the limit and neither confirm nor contradict the expansions. The ratios $c_n/g_n\approx0.952034$, $0.956472$ at $n=12$, $13$ exceed $1-\frac1{2n}-\frac1{n^2}$ (Part~II, \eqref{cgr:fc:eq:ratio-main}) by $6.45\cdot10^{-4}$ and $8.50\cdot10^{-4}$; $n^3$ times the excess is $1.11$ and $1.87$, so at these two orders the excess is of the size of an $O(n^{-3})$ remainder but has not settled to a constant.
\end{writenote}

\subsection{The reproducible companion}\label{cgr:ld:sub:companion}
The newly written standard-library Python companion verifies
finite instances of the fixed-endpoint leaf completion formula
\eqref{cgr:ld:eq:leaf-completions}, enumerates matching and graph
isomorphism classes at stated small orders, and checks explicit
independent leaf-move constructions. In each eligible example it
checks the graph, every dihedral orbit size, and disjointness of
the $2^k$ orbits. It distinguishes geometric symmetries from graph
automorphisms rather than conflating the two counts.

Further exact checks reconcile the published graph counts and the
component product, and verify the inverse coefficient algebra.
The machine-readable evidence file records all finite ranges and
examples. A separate validator recomputes the evidence and rejects
semantic alterations; tests run both normally and with Python
optimization enabled, so no mathematical checks depend on
removable \texttt{assert} statements.

The finite tests check implementation, normalization and examples.
They cannot establish a limiting Poisson law, rule out every
large-$n$ exceptional representation, prove the recovery lemma
for all cores, or establish the quantitative asymptotic rates.
Those tasks are handled by the mathematical arguments and cited
structural inputs. Small forcing-coupling witnesses likewise do not
replace the all-size argument of Section~\ref{cgr:ld:sec:quantitative}.
\begin{writenote}[The companion here, 5~October 2026]
Shipped as \file{code/150-leading-companion-*.py} and \file{data/150-leading-companion-*} (the README maps the delivered names). At intake its evidence was regenerated byte for byte and validated on a copy; its unit tests and the release tests need POSIX file calls and were not run on Windows (README, ``Rerunning the checks'').
\end{writenote}

\section{Limitations and further work}\label{cgr:ld:sec:limits}
The quantified sandwich proves a relative $O(n^{-1})$ error but
still does not determine the individual first correction coefficient.
Several extensions would require additional work.
\begin{enumerate}
\item A first individual coefficient requires analysis of the
order-$1/n$ exceptional configurations, including long leaves,
shared parents, overlapping local decorations and nonlocal splits,
and their actual representation fibers. The connectivity correction
alone does not identify that coefficient.
\item An all-orders expansion cannot be inferred by freely
expanding the decorated-core lower bound. Nongeneric cores and
other decompositions must be controlled at the requested order;
the current $O(n^{-1})$ exceptional mass cannot simply be dropped.
\item Effective graph-count constants and onsets would make
\eqref{cgr:ld:eq:quant-inverse} a numerical threshold certificate.
The explicit $200/n$ bound applies only to the local matching
patterns, not the complete graph-count argument.
\item The final 2026 journal text should be reconciled with the
2024 sources before any stronger priority statement. The optional
labeled consequence deserves its own source and sequence audit.
\end{enumerate}
The central model-specific conclusions established here are the
reciprocal-fiber upper bound, the canonically recoverable lower
construction, their exact matching normalization, and the
consequences forced by saturation.
\begin{writenote}[Status of these items, 5~October 2026]
Item~1 is answered by Part~II's Theorems~\ref{cgr:fc:thm:main} and~\ref{cgr:fc:thm:transfer-summary}, by another route than the one sketched here: instead of analysing the exceptional configurations listed, Part~II expands the local pattern counts uniformly (Theorem~\ref{cgr:fc:thm:local}), controls the nonlocal decompositions by its three-chord interval lemma, and transfers the prime correction to graph classes through canonical split trees. Item~2 is re-scoped: Part~II's conditional transfer (Subsection~\ref{cgr:fc:sub:conditional}) controls the nongeneric cores and the graph splits at every fixed order and reduces the graph expansion to the expansion of $p_n^{\mathrm a}/h_n$, which is known only through its first correction; the next prime coefficient $\beta$ is open (Section~\ref{cgr:sec:further}, item~\ref{cgr:q:beta}). Items~3 and~4 remain open (items~\ref{cgr:q:effective}, \ref{cgr:q:sources} and~\ref{cgr:q:labeled}).
\end{writenote}

\section{Dependency map and reproducible release}\label{cgr:ld:app:dependency}
\begin{writenote}
Report~150's Appendix~A, printed as a numbered section of Part~I so that Part~II's sections follow it.
\end{writenote}
The logical dependencies are as follows.
\begin{enumerate}
\item The published matching decomposition and local Poisson
lemmas give prime probability $e^{-3}$; prime uniqueness and the
geometric automorphism action connect this to asymmetric prime
graph counts.
\item The direct rotation and two-reflection estimates control
all geometric symmetries, including symmetries induced after
bounded deletions.
\item Long-leaf and shared-parent bounds plus independent leaf
moves bound connected reciprocal fibers from above.
\item Canonical recovery of the three decorations makes the
lower count injective, and its normalized exponential sum matches
the upper constant.
\item The component convolution transfers the equivalent to all
graphs and isolates the connected/all correction. Saturation
gives the core and automorphism laws. A direct local-pattern
forcing coupling, an elementary Poisson approximation bound and a
logarithmic cutoff quantify these rates. Monotone smooth inversion
gives the threshold bracket.
\end{enumerate}
The source archive contains the report PDF, editable LaTeX,
new companion code and evidence, a source-provenance record,
checksums, deterministic build scripts, and release tests. It
excludes downloaded papers, private working notes, prior reports,
and build intermediates. The README gives the exact commands,
required toolchain, finite computation ranges and output-safety
rules. PDF byte reproducibility is qualified to the tested
pdfTeX toolchain; the mathematical contents are not dependent on
that byte identity. The checksum manifest establishes package
consistency, not authenticity against replacement of both a
payload and its manifest.
\begin{writenote}[The release here, 5~October 2026]
The PDF and the checksum manifest \file{SHA256SUMS} are not shipped (repository policy; the manifest verified 17/17 at intake); the editable source is this Part; the companion, provenance record and release scripts are shipped under prefixed names (README).
\end{writenote}

\clearpage
\part{First corrections for unlabeled circle graphs}\label{cgr:fc:part}
\partsource{Bundle Report~152, archive \file{Circle_Graph_First_Corrections_and_Refined_Inverses_Source.zip}, delivered as \file{Report152.tex} (1200 lines, 19 pages). Delivered title block ``First corrections for unlabeled circle graphs / Connected and all graphs in OEIS A156808 and A156809 / Report 152 / 3 October 2026''. Printed as delivered, with \textbf{[write]} notes. Delivered Section~$k$ is Section~$k+13$ here, and its statement and equation numbers $k.j$ are $(k+13).j$. Section~\ref{cgr:sec:further}, which closes the report, is added in the write.}
\begin{center}
\small
\begin{tabular}{@{}>{\raggedright\arraybackslash}p{0.46\textwidth} >{\raggedright\arraybackslash}p{0.46\textwidth}@{}}
\toprule
Reading conventions of Part~II & Elsewhere in the report\\
\midrule
$c_r$ the coefficients of $H_M(u)$ (Section~\ref{cgr:fc:sec:polymer}) & $c_n=\mathrm{A156808}(n)$ everywhere else\\
$q_r=\prod_{j<r}(M-2j-1)^{-1}$ & Part~I: $\rho_r$; Part~I's $q_s$ is $m_sm_{n-s}/m_n$\\
$x,y,z$ complex weights; $y$ also the threshold input & Part~I: $x$ the threshold input\\
$\alpha=-5$, $\beta$ the prime coefficients; $\gamma=\alpha+\frac54$, $\eta=\gamma+\frac12$ & Part~I: $\alpha,\beta$ pattern indices; $\eta$ bracket errors\\
$B=\frac{3+\log2}2$ and $D=\log(2u)$ in the inverse & Part~I: $B=\frac32(1+\log2)$, $A=\log(2u)$\\
$b_j$, $b_j(t)$ rooted counts; $r_s$ rooted classes; $\rho=2^{-t}$, $\sigma=6^{-t}$ & Part~I: $\rho_r$ a pattern probability\\
$C_n^{\mathrm{lab}}$, $G_n^{\mathrm{lab}}$ & Part~I: $c_n^{\mathrm{lab}}$, $g_n^{\mathrm{lab}}$\\
\bottomrule
\end{tabular}
\end{center}
\begin{partabstract}
Let $c_n$ and $g_n$ count connected and arbitrary unlabeled circle
graphs on $n$ vertices, and put $h_n=(2n-1)!!/(4n)$. We prove
\[
 c_n=e^{-3/2}h_n\left(1-\frac{15}{4n}+O(n^{-2})\right),\qquad
 g_n=e^{-3/2}h_n\left(1-\frac{13}{4n}+O(n^{-2})\right).
\]
A uniform complex generating-function expansion for three local
matching patterns gives prime probability
$e^{-3}(1-5/n+O(n^{-2}))$. A direct three-chord interval lemma
controls the nonlocal remainder. Canonical graph split trees then
transfer the prime coefficient to graph-isomorphism classes, with
explicit rooted-branch counts and globally controlled tails.
We also obtain $c_n/g_n=1-1/(2n)-1/n^2+O(n^{-3})$ and a refined
Lambert-$W$ inverse with ceiling uncertainty $O(1/(u^2\log u))$.
An every-fixed-order transfer theorem is conditional on the
corresponding prime expansion; no unconditional higher prime
coefficient is claimed. Exact finite diagnostics and deterministic
PDF and source-archive builders accompany the proof.
\end{partabstract}
\noindent\textbf{Source and scope.} The primary circle-graph source
basis is arXiv:2402.06394v1 (2024), a complete author-hosted 2024
copy, and the structural references below. The 2026 journal version
was identified bibliographically; its final text was not inspected.
The first corrections established here were not located in the
bounded source review. This is not an exhaustive priority claim.
The remainders and inverse onsets are asymptotic existence bounds,
not effective finite-input certificates. No external publication,
author contact, or public repository change is part of this report.
\begin{writenote}[Reading Part~II against Part~I, 5~October 2026]
``The earlier Report~150'' (Subsection~\ref{cgr:fc:sub:priority}) is Part~I. Part~II's $c_n$, $g_n$, $m_n$, $h_n$, $\mathcal M_n$, $S_n$, $R_n(G)$, $p_n^{\mathrm a}$, $X,Y,Z$ and $M=2n$ are Part~I's. Part~II cites Part~I only in that priority paragraph: it re-proves Part~I's symmetry and coarse bounds (Subsection~\ref{cgr:fc:sub:symmetry}) and proves its local expansion without the qualitative Poisson limit Part~I starts from. Its Theorem~\ref{cgr:fc:thm:main} supersedes the count statements of Part~I's Theorem~\ref{cgr:ld:thm:quantitative}, and its inverse refines Part~I's Section~\ref{cgr:ld:sec:inverse}.
\end{writenote}

\section{Results and conventions}\label{cgr:fc:sec:results}
A circle graph is the intersection graph of chords with distinct
endpoints on a circle. All graphs are finite, simple and undirected;
``unlabeled'' means one graph-isomorphism class, counted once.
Write
\[
 c_n=\mathrm{A156808}(n),\quad g_n=\mathrm{A156809}(n),\quad
 m_n=(2n-1)!!=\frac{(2n)!}{2^n n!},\quad h_n=\frac{m_n}{4n}.
\]
The two OEIS offsets are one. We use $m_0=g_0=1$ for the empty
matching and graph. An indexed matching has its $2n$ endpoints in
fixed cyclic order, but no independent labels on its chords. The
dihedral group on the endpoints has order $4n$. Geometric symmetry
of a matching and abstract symmetry of its graph are different
notions; we identify them only for the prime class where justified.

A graph split is a partition $A\sqcup B$ with $|A|,|B|\ge2$
whose crossing edges form $A'\times B'$ for subsets of the two
sides. A nondegenerate split-prime graph has no such split;
cliques and stars are treated as degenerate labels in split trees.
Let $p_n^{\mathrm a}$ count abstractly asymmetric prime circle
graphs. All assertions about primes concern sufficiently large
orders, so exceptional small conventions play no role.

\begin{theorem}[Individual first corrections]\label{cgr:fc:thm:main}
As $n\to\infty$,
\begin{align}
 p_n^{\mathrm a}&=e^{-3}h_n\left(1-\frac5n+O(n^{-2})\right),
 \label{cgr:fc:eq:prime-main}\\
 c_n&=e^{-3/2}h_n\left(1-\frac{15}{4n}+O(n^{-2})\right),
 \label{cgr:fc:eq:c-main}\\
 g_n&=e^{-3/2}h_n\left(1-\frac{13}{4n}+O(n^{-2})\right).
 \label{cgr:fc:eq:g-main}
\end{align}
Equivalently, with $Q_n=e^{-3/2}(2n/e)^n/(2\sqrt2\,n)$,
\[
 c_n=Q_n\left(1-\frac{91}{24n}+O(n^{-2})\right),\qquad
 g_n=Q_n\left(1-\frac{79}{24n}+O(n^{-2})\right).
\]
The latter coefficients include the $-1/(24n)$ correction in
$m_n$; they are not the coefficients in the $h_n$ normalization.
\end{theorem}
\begin{writenote}[Relation to Part~I, 5~October 2026]
With $h_n=Q_n(1-\frac1{24n}+O(n^{-2}))$, the two normalizations differ by $-\frac1{24}$ in the first coefficient. Both expansions imply Part~I's Theorems~\ref{cgr:ld:thm:main} and~\ref{cgr:ld:thm:quantitative} for the counts.
\end{writenote}

\begin{theorem}[Connectivity and conditional higher transfer]
\label{cgr:fc:thm:transfer-summary}
Unconditionally,
\begin{equation}\label{cgr:fc:eq:ratio-main}
 \frac{c_n}{g_n}=1-\frac1{2n}-\frac1{n^2}+O(n^{-3}).
\end{equation}
For each fixed $q\ge0$, an expansion of $p_n^{\mathrm a}/h_n$
through $n^{-q}$ with remainder $O(n^{-q-1})$ transfers to such
expansions of both $c_n/h_n$ and $g_n/h_n$. This is a conditional
transfer theorem, not an unconditional full expansion. In
particular, if the additional hypothesis
\[
 p_n^{\mathrm a}=e^{-3}h_n\left(1-\frac5n+
                  \frac\beta{n^2}+O(n^{-3})\right)
\]
holds for some $\beta$, then the second relative coefficients are
$\beta-277/32$ for $c_n$ and $\beta-297/32$ for $g_n$.
No value of $\beta$ is proved or conjectured here.
\end{theorem}

The main two proof steps are different. The local-pattern analysis
counts indexed matchings. The graph transfer counts unlabeled
isomorphism classes via their canonical split trees. In particular,
we never assume that a graph split of a specified size is an
interval split of that same size in every chord representation.

\section{Structural source inputs and a coarse bound}\label{cgr:fc:sec:inputs}
Bassino, Bouvel, F\'eray, Gerin and Pierrot
\cite[Definition 5.3, Proposition 5.5]{BBFGP24} define a matching
decomposition by four successive circular intervals, with every
chord inside one interval or joining opposite intervals. Empty
intervals are allowed. In their anchored $k$-decomposition,
$C_1$ contains a fixed endpoint and the $k$ chords lie in
$C_2\cup C_4$. Their structural equivalence says that an
indecomposable matching has a split-prime intersection graph and
conversely. We use their Lemma 5.9 to identify zero local defects
with the absence of the two endpoint decomposition sizes, and
Lemma 5.6's exact count
\begin{equation}\label{cgr:fc:eq:source-decomp}
 d_n^k=(n-k)m_{k+1}m_{n-k+1}
\end{equation}
for pairs consisting of a matching and one chosen $k$-decomposition.
It upper-bounds the number of distinct matchings admitting such a
decomposition; their Lemma 5.7 uses it to bound the summed
probability. Their Proposition 5.12, attributed there to
Gabor, Supowit and Hsu, gives uniqueness of a prime representation
up to rotation and reflection for order at least five.

Klav\'ik and Zeman \cite[Section 3, Lemmas 7--8 and their
proofs]{KZ15} give closure of circle graphs under splitting and
inverse split composition. They identify graph automorphisms with
canonical split-tree automorphisms and describe prime graph
automorphisms as rotations or reflections of the unique circle
representation. The action on chords is faithful in prime order at
least five. Indeed, a rotation fixing every chord must be a
half-turn with all chords diameters, yielding a clique. A
reflection fixing every chord gives noncrossing chords perpendicular
to its axis, with at most one additional axis chord; its graph is
a star with possible isolates. Neither graph is prime at this order.
Thus prime abstract asymmetry is equivalent to geometric asymmetry.

Gioan and Paul \cite[Section 2.2, Theorem 2.9]{GP12} give the
unique reduced graph-labelled split tree of a connected graph.
Its internal nodes have degree at least three, with prime, clique
or star labels. Adjacent clique labels and the reducible
center-to-leaf star joins are forbidden. Their Lemma 2.8 says each
internal tree edge induces a graph split. Corollary 2.10 says every
graph split is induced by an edge after at most one split of a
degenerate node. Corollary 2.6 identifies labels as induced
subgraphs; Lemma 2.3 supplies connectedness of the accessibility
graph when all labels are connected. These precise source inputs
are what the graph-level proof needs.

\subsection{Symmetries are negligible to every fixed order}\label{cgr:fc:sub:symmetry}
Let $S_n$ be the number of indexed matchings fixed by a nonidentity
rotation or reflection. A rotation of order $d\ge2$ has
$k=2n/d$ endpoint cycles. Its invariant matchings number
\[
 k![z^k]\exp\left(\one_{2\mid d}z+\frac d2z^2\right).
\]
Two different cycles are paired with one of $d$ offsets; an
internally matched cycle must have even order and pair opposite
points. Evaluating the nonnegative series at $z=\sqrt{k/d}$ and
using $k!\le k^k$ bounds this count by
\[
 (kd)^{k/2}\exp(k/2+\sqrt{k/d})
 \le (2en)^{n/2}\exp(\sqrt n).
\]
A gap-axis reflection has the same count as $d=2,k=n$. An
endpoint-axis reflection forces its two fixed endpoints to match
together, leaving the case $d=2,k=n-1$. Therefore
\begin{equation}\label{cgr:fc:eq:symmetry}
 S_n\le4n(2en)^{n/2}e^{\sqrt n}=o(m_n n^{-A})
 \quad\hbox{for every fixed }A>0.
\end{equation}
The last claim follows from $\log m_n=n\log(2n)-n+O(1)$.

If $R_n(G)$ is the number of indexed representations of $G$,
then grouping matchings by their graph gives exactly
\[
 g_n=\sum_{D\in\mathcal M_n}\frac1{R_n(G_D)}.
\]
For a geometrically asymmetric $D$, its orbit has $4n$ members,
so $R_n(G_D)\ge4n$; on the other matchings use $1/R_n\le1$.
Consequently
\begin{equation}\label{cgr:fc:eq:coarse}
 c_n\le g_n\le h_n+S_n=O(h_n),\qquad
 g_s\le C\frac{m_s}{s}\quad(s\ge1)
\end{equation}
for a constant enlarged over finitely many small orders. This
short argument makes the new transfer proof independent of any
unproved strengthening of earlier leading estimates.
\begin{writenote}[Restated from Part~I, 5~October 2026]
\eqref{cgr:fc:eq:symmetry} is Part~I's Lemma~\ref{cgr:ld:lem:symmetry} with \eqref{cgr:ld:eq:symbound}, by the same argument; the fibre identity and \eqref{cgr:fc:eq:coarse} are Part~I's \eqref{cgr:ld:eq:fiber}, \eqref{cgr:ld:eq:coarse-g} and \eqref{cgr:ld:eq:uniform-g}, with $c_n\le g_n$ added. The method is that of \cite[Proposition 5.14]{BBFGP24}; see Remark~\ref{cgr:ld:rem:symmetry-credit}. The faithfulness argument of the preceding paragraphs is Part~I's, before \eqref{cgr:ld:eq:prime-asym}.
\end{writenote}

\subsection{Bounded source priority}\label{cgr:fc:sub:priority}
The complete 2024 arXiv v1 and the complete author-hosted 2024 copy
were inspected in the preceding source review; the relevant text
agrees. The journal metadata is \emph{Electronic Journal of
Probability} \textbf{31} (2026), article 110,
\href{https://doi.org/10.1214/26-EJP1570}{doi:10.1214/26-EJP1570}.
Its final text was not inspected. The earlier Report 150 established
the leading equivalent and relative $O(1/n)$ bounds; this report
identifies the individual first coefficients with stronger errors.
Neither that comparison nor the bounded literature search
certifies universal novelty. The source archive records the
reviewed versions, source hashes, and proof dependencies.
\begin{writenote}[Sources, 5~October 2026]
As in Part~I (notes in Subsections~\ref{cgr:ld:sub:poisson} and~\ref{cgr:ld:sub:inspected}): arXiv:2402.06394v1 is the only arXiv version as of 5~October 2026 and was read at the write; the journal text was not read. The write also read Gioan and Paul's statements cited in this section in arXiv:0810.1823v2, where they carry the same numbers; the published version was not compared.
\end{writenote}

\section{Local atoms and the exact polymer representation}
\label{cgr:fc:sec:polymer}
Put $M=2n$, and work on the cycle with vertex set
$\mathbb Z/M\mathbb Z$. An $X$ atom requires a chord between
consecutive endpoints. A $Y$ atom requires a chord at cyclic
distance two. A $Z$ atom requires one of the two matchings between
two disjoint consecutive endpoint pairs. Count each such specified
two-chord configuration once. For $M\ge6$ these give $M$ atoms of
each singleton type and $M(M-3)$ atoms of type $Z$. Let $X,Y,Z$
be the counts of atoms present in a uniform matching. The letters
$x,y,z$ below are complex weights, not counting variables for
vertices.

\begin{theorem}[Uniform local expansion]\label{cgr:fc:thm:local}
Uniformly for $(x,y,z)$ in every fixed compact subset of
$\mathbb C^3$,
\begin{equation}\label{cgr:fc:eq:local-exp}
 \EE[(1+x)^X(1+y)^Y(1+z)^Z]
 =e^{x+y+z}\left(1+\frac{F(x,y,z)}M+O(M^{-2})\right),
\end{equation}
where
\begin{equation}\label{cgr:fc:eq:F}
 F=x+y+z-\tfrac12x^2-\tfrac12y^2-z^2-2xy-xz-yz+y^2z.
\end{equation}
In particular,
\begin{equation}\label{cgr:fc:eq:localzero}
 \PP(X=Y=Z=0)=e^{-3}\left(1-\frac{10}M+O(M^{-2})\right).
\end{equation}
\end{theorem}
\begin{writenote}[Monte Carlo, uncertified, 5~October 2026]
At intake, $\PP(X=Y=Z=0)$ was estimated on 200{,}000 uniform matchings for each of $M=100$, $200$, $400$ (NumPy, seed 20261005; intake record, not shipped). The rescaled estimates of the coefficient of $1/M$, namely $M(\widehat P/e^{-3}-1)$, are $-9.88\pm0.93$, $-9.83\pm1.9$ and $-4.4\pm3.9$, with one binomial standard error. The first two agree with $-10$ and exclude $0$ by about $10$ and $5$ standard errors; at $M=400$ the estimate lies within $1.5$ standard errors of both $-10$ and $0$. This is evidence for a proved statement, and it says nothing about the second-order term.
\end{writenote}

\subsection{Exact partition polynomial}
An atom set $A$ is compatible if its required chords form a
matching. Write $r(A)$ for their number of distinct chords and
$t(A)$ for the product of the atom weights. Define
\begin{equation}\label{cgr:fc:eq:partition}
 H_M(u)=\sum_{A\text{ compatible}}t(A)u^{r(A)}
       =\sum_{r=0}^{n}c_r u^r.
\end{equation}
Expansion of the product of atom indicators is exact and gives
\begin{equation}\label{cgr:fc:eq:matching-exact}
 \EE[(1+x)^X(1+y)^Y(1+z)^Z]=\sum_{r=0}^{n}c_rq_r,
 \qquad q_r=\prod_{j=0}^{r-1}(M-2j-1)^{-1}.
\end{equation}
Here $q_r$ is the probability of $r$ specified disjoint chords.
\begin{writenote}
In this section and the next, $c_r$ is a coefficient of $H_M(u)$, not $\mathrm{A156808}(r)$.
\end{writenote}

Connect two atoms when their endpoint supports meet. Inside a
compatible set this means they share a whole chord. A polymer is
a nonempty compatible atom set connected under this relation. Its
support is denoted $S_P$, size $r(P)$ counts its distinct chords,
and activity is $w_P=t(P)u^{r(P)}$. Compatible atom sets have unique
connected-component decompositions into endpoint-disjoint polymers;
conversely such a family reconstructs the atom set. Hence
$H_M$ is exactly the hard-core polymer partition polynomial,
where overlapping supports, including identical polymers, conflict.

\subsection{Uniform counting and absolute tails}\label{cgr:fc:sec:polymer-bounds}
In the graph of possible chords, join two chords when together
they form a $Z$ atom. Each chord has at most four neighbors: choose
one cycle neighbor of each endpoint. A connected chord set of size
$r\ge2$ with a specified root chord has at most $4^{2r-2}$
possibilities. To see this, choose a canonical spanning tree and
its depth-first traversal of length $2r-2$; each step has at most
four choices and the visited set determines the chord set.
There are at most $2r$ possible $Z$ edges on that set, and at most
$r$ possible singleton atoms. Thus at most $2^{3r}$ atom subsets
occur on it. Bounded complex weights can be absorbed into $C^r$.
For a constant depending only on the fixed weight compact set:
\begin{itemize}
\item There are at most $M^2C^r$ polymers of size $r\ge2$
\item At most $MC^r$ of them contain any specified endpoint,
by rooting at its unique incident chord
\item There are exactly $2M$ size-one polymers
\item At most $2MC^r$ polymers of size $r$ contain a singleton
atom, by rooting at that singleton chord
\end{itemize}
For every fixed $S$ and nonnegative integer $k$, these estimates
and geometric summation give, uniformly for $|u|\le S/M$,
\begin{equation}\label{cgr:fc:eq:polymer-bounds}
 \sum_P r(P)^k|w_P|=O(1),\qquad
 \sup_v\sum_{P:v\in S_P}r(P)^k|w_P|=O(M^{-1}).
\end{equation}
The total activity of size at least four is $O(M^{-2})$, with any
fixed polynomial size weight. So is that of size at least three
with a singleton atom. In addition, dropping endpoint restrictions
in the positive-activity polymer sum yields
\begin{equation}\label{cgr:fc:eq:absolute-coeff}
 \sum_r|c_r|(S/M)^r
 \le\sum_{A\text{ compatible}}|t(A)|(S/M)^{r(A)}
 \le\exp\left(\sum_P|t(P)|(S/M)^{r(P)}\right)=O(1).
\end{equation}
The fact that this holds for every fixed $S$ is essential. Fixed
factorial moments alone would not justify the later summation over
all $r$ or evaluation at the zero-defect weights.

\section{The first hard-core correction}\label{cgr:fc:sec:hardcore}
For fixed $u$, put
\[
 W=\sum_Pw_P,\qquad B=\sum_{P\sim Q}w_Pw_Q,
\]
where the latter sum is ordered and includes $P=Q$. Then
\begin{equation}\label{cgr:fc:eq:hardcore}
 H_M(u)=e^W(1-B/2)+O(M^{-2})\qquad(|u|\le S/M).
\end{equation}
Here is an elementary remainder proof valid for complex activities.
Express the partition function as the sum over ordered polymer
$k$-tuples, divided by $k!$, with the indicator of no conflicting
index pairs. If $N$ is the number of such pairs in a tuple, then
\[
 \left|\one_{\{N=0\}}-(1-N)\right|\le\binom N2.
\]
The terms $1$ and $-N$ sum to $e^W$ and $-e^WB/2$.
Two distinct conflict pairs use either four or three indices.
With absolute activities the resulting remainder is bounded by
\begin{equation}\label{cgr:fc:eq:two-conflicts}
 e^{W_*}(B_*^2/8+C_*/2),
\end{equation}
where
\[
 W_*=\sum_P|w_P|,\quad B_* =\sum_{P\sim Q}|w_Pw_Q|,
 \quad C_* =\sum_P|w_P|\left(\sum_{Q\sim P}|w_Q|\right)^2.
\]
By endpoint incidence in \eqref{cgr:fc:eq:polymer-bounds},
$\sum_{Q\sim P}|w_Q|=O(r(P)/M)$. Therefore $W_*=O(1)$,
$B_*=O(M^{-1})$ and $C_*=O(M^{-2})$. This proves
\eqref{cgr:fc:eq:hardcore} without appealing to a formal infinite cluster
series or assuming positivity of the original weights.

\subsection{All relevant connected polymers}
Set $u=s/M$ for bounded complex $s$. Size-one polymers contribute
$s(x+y)$; single-$Z$ polymers contribute
$zs^2-3zs^2/M$. The other two-chord polymers are:
\begin{center}
\begin{tabular}{@{}lll@{}}
\toprule
Atom types & Number & Reason\\
\midrule
$XZ$ & $M$ & The partner of $\{i,i+1\}$ is $\{i-1,i+2\}$\\
$YZ$ & $3M$ & Three partners of each $\{i,i+2\}$\\
$YYZ$ & $M$ & $\{i,i+2\},\{i+1,i+3\}$\\
\bottomrule
\end{tabular}
\end{center}
For the $Y$ chord the three partners are
$\{i-1,i+1\},\{i-1,i+3\},\{i+1,i+3\}$.
Distinct singleton atoms cannot share a chord. This list is
therefore exhaustive on two chords. In particular, the two $YZ$
subpolymers in a $YYZ$ configuration are different from its full
three-atom polymer; all must be counted. Their combined extra
activity is $s^2(xz+3yz+y^2z)/M$.

The $Z$-adjacency graph of any compatible matching has maximum
degree two. Two compatible neighbors of a chord must use the two
different cycle neighbors of each endpoint. A cycle of such
adjacencies therefore has endpoint support closed under neighbors
of the original cycle, and hence uses all $M$ endpoints. Thus
three-chord $Z$ cycles are impossible for $M\ge8$.
The three-chord polymers without singleton atoms are paths of two
$Z$ edges. Their number is exactly $M(M-5)$. Indeed the common
chord has distance at least three in both directions, giving
$M(M-5)/2$ choices; its four neighboring endpoints have two
compatible pairings. The common chord is unique. This contributes
$z^2s^3/M+O(M^{-2})$. Section~\ref{cgr:fc:sec:polymer-bounds} controls
every remaining polymer. Hence
\begin{align}
 W&=A(s)+W_1(s)/M+O(M^{-2}),\label{cgr:fc:eq:W}\\
 A(s)&=s(x+y)+s^2z,\notag\\
 W_1(s)&=s^2(-3z+xz+3yz+y^2z)+s^3z^2.\notag
\end{align}
Both the $YYZ$ term and the shared-chord $ZZ$ path contribute at
first order. Truncating original-atom clusters at pairs would miss
one of them.

\subsection{Endpoint exclusions}
A polymer is elementary if it contains one original atom.
Non-elementary polymers have total absolute activity $O(M^{-1})$,
even with a polynomial size weight. Their contribution to $B$ is
$O(M^{-2})$ by endpoint incidence. Elementary supports have the
following ordered overlap counts, including identical atoms:
\begin{equation}\label{cgr:fc:eq:overlap-table}
\begin{array}{c|rrrrrr}
 \text{type order}&XX&YY&XY&XZ&YZ&ZZ\\\hline
 \text{count}&3M&3M&4M&6M^2+O(M)&8M^2+O(M)&12M^3+O(M^2).
\end{array}
\end{equation}
The first three follow by fixing the first singleton and choosing
one of three, three or four intersecting singleton supports.
Every endpoint lies in $4(M-3)$ $Z$ supports. An adjacent pair
lies in $2M-4$ of them, so an $X$ support meets $6M-20$ $Z$
atoms. A distance-two pair lies in six $Z$ supports, so a $Y$
support meets $8M-30$ of them.

For the last leading constant, inclusion--exclusion over shared
endpoints starts with $M[4(M-3)]^2=16M^3+O(M^2)$.
Subtracting shared adjacent pairs removes
$M(2M-4)^2=4M^3+O(M^2)$. Nonadjacent pairs are in at most a
fixed number of $Z$ supports, so their total is $O(M^2)$.
Any triple or quadruple is also in only a fixed number of $Z$
supports, and only $O(M^2)$ such sets appear in any support.
Thus the remaining inclusion--exclusion terms are $O(M^2)$,
proving the $12M^3$ leading term.
Consequently
\begin{align}
 B&=B_1(s)/M+O(M^{-2}),\label{cgr:fc:eq:B}\\
 B_1(s)&=s^2(3x^2+3y^2+8xy)
       +s^3(12xz+16yz)+12s^4z^2.\notag
\end{align}
Combining \eqref{cgr:fc:eq:hardcore}, \eqref{cgr:fc:eq:W} and \eqref{cgr:fc:eq:B},
\begin{equation}\label{cgr:fc:eq:H-exp}
 H_M(s/M)=e^{A(s)}
 \left(1+\frac{W_1(s)-B_1(s)/2}M\right)+O(M^{-2}),
\end{equation}
uniformly on every fixed complex compact set of the weights and $s$.

\section{Exact matching probabilities and the prime correction}
\label{cgr:fc:sec:matching}
\subsection{A global remainder through the full matching range}
There is an absolute constant $C$ such that, for $0\le r\le M/2$,
\begin{equation}\label{cgr:fc:eq:q-global}
 q_r=M^{-r}\left(1+\frac{r^2}M+R_{r,M}\right),\qquad
 |R_{r,M}|\le\frac{Cr^4e^r}{M^2},\quad R_{0,M}=0.
\end{equation}
To prove this, put
\[
 a_{r,M}=\log(q_rM^r)=\sum_{j=0}^{r-1}
           -\log\left(1-\frac{2j+1}M\right).
\]
Convexity and the midpoint-integral inequality give
\[
 a_{r,M}\le\frac M2\int_0^{2r/M}-\log(1-v)\,dv\le r,
\]
including $r=M/2$ by the integrable endpoint limit. For $r\le M/4$,
Taylor's formula for $-\log(1-v)$ gives
\[
 0\le a_{r,M}-r^2/M\le Cr^3/M^2,\qquad a_{r,M}\le2r^2/M.
\]
Using $|e^a-1-a|\le a^2e^a/2$ proves \eqref{cgr:fc:eq:q-global} in
this range. For $r>M/4$, use $e^a\le e^r$ and
$r^4/M^2\ge1$ for large $M$, increasing $C$ for the finite
remaining cases. Thus the estimate is global, not merely a
fixed-$r$ approximation.

Insert this into \eqref{cgr:fc:eq:matching-exact}. Since
$r^4e^r\le C'e^{2r}$, \eqref{cgr:fc:eq:absolute-coeff} with $S=e^2$
makes the total error $O(M^{-2})$. Writing $\mathcal D=s\,d/ds$,
\begin{equation}\label{cgr:fc:eq:q-summed}
 \sum_rc_rq_r=H_M(1/M)+\frac1M
       \left.\mathcal D^2H_M(s/M)\right|_{s=1}+O(M^{-2}).
\end{equation}
Uniform complex analyticity of \eqref{cgr:fc:eq:H-exp} near $s=1$
permits differentiation of its leading approximation by Cauchy's
estimate. Therefore
\[
 \left.\mathcal D^2H_M(s/M)\right|_{s=1}
 =e^{x+y+z}\bigl((x+y+2z)^2+x+y+4z\bigr)+O(M^{-1}).
\]
The relative coefficient in \eqref{cgr:fc:eq:q-summed} is
\[
 W_1(1)-B_1(1)/2+(x+y+2z)^2+x+y+4z=F(x,y,z).
\]
This proves Theorem~\ref{cgr:fc:thm:local}; at $x=y=z=-1$, $F=-10$.
For example, the univariate consequence is
\[
 \EE(1+t)^{X+Y+Z}
 =e^{3t}\left(1+\frac{3t-6t^2+t^3}M+O(M^{-2})\right).
\]
This identity follows from uniform analyticity; fixed cumulants
were not used as a substitute for the global bounds.

\subsection{A direct three-chord interval lemma}
\begin{lemma}\label{cgr:fc:lem:three}
If all endpoints of three chords occupy the union of two circular
intervals containing no endpoints of other chords, then those
three chords contain an $X$, $Y$ or $Z$ atom.
\end{lemma}
\begin{proof}
Swap intervals so their endpoint counts are $0+6$, $1+5$,
$2+4$ or $3+3$. An internal chord of within-interval distance
one or two already gives $X$ or $Y$, so suppose none exists.
\begin{itemize}
\item In $0+6$, positions $2,1,0$ must match respectively
$5,4,3$. Two parallel consecutive chords give $Z$
\item In $1+5$, the only two disjoint internal chords of length
at least three are $\{0,3\}$ and $\{1,4\}$; these give $Z$
\item In $2+4$, both endpoints of the shorter interval match into
the longer one. Its remaining internal chord must be $\{0,3\}$;
positions $1,2$ match to the adjacent shorter-interval pair,
giving $Z$
\item In $3+3$, all chords join the intervals. The two adjacent
pairs of the first interval have distinct image pairs; only one
pair of the second interval is nonadjacent. At least one image
pair is adjacent, giving $Z$
\end{itemize}
Extra adjacency across an empty separating interval can only add
local atoms, so the proof covers that case too.
\end{proof}

A $3$- or $(n-3)$-decomposition has one side satisfying the
lemma, while source Lemma 5.9 excludes $2$ and $n-2$ when
$X=Y=Z=0$. Thus a zero-defect decomposable matching must have a
$k$-decomposition with $4\le k\le n-4$. Set
$e_n^k=m_{k+1}m_{n-k+1}$; it is symmetric and decreases towards
the midpoint because
\[
 \frac{e_n^{k+1}}{e_n^k}=\frac{2k+3}{2n-2k+1}.
\]
The source bound \eqref{cgr:fc:eq:source-decomp} and a union bound give
\begin{align*}
 \sum_{k=4}^{n-4}\frac{d_n^k}{m_n}
 &\le n\left(2\frac{e_n^4}{m_n}+n\frac{e_n^5}{m_n}\right)\\
 &=n\bigl(O(n^{-3})+nO(n^{-4})\bigr)=O(n^{-2}).
\end{align*}
This is a statement about decompositions of the actual matching,
with no graph-split-size assumption. Consequently
\begin{equation}\label{cgr:fc:eq:prime-prob}
 \PP(D_n\text{ indecomposable})
 =\PP(X=Y=Z=0)+O(n^{-2})
 =e^{-3}\left(1-\frac5n+O(n^{-2})\right).
\end{equation}
The difference from the zero-defect probability is nonpositive.
Removing the at most $S_n$ symmetric matchings and dividing each
remaining prime orbit by exactly $4n$ proves
\eqref{cgr:fc:eq:prime-main}, by prime uniqueness and
\eqref{cgr:fc:eq:symmetry}.
\begin{writenote}[Credit and relation to Part~I, 5~October 2026]
The ratio $e_n^{k+1}/e_n^k$, the monotonicity towards the midpoint and the bound of the middle terms by their extreme values are those of the source's proof of its Lemma~5.7 (with the same notation $e_n^k$). Part~II's sum starts at $k=4$, because Lemma~\ref{cgr:fc:lem:three} excludes $k=3$ and $k=n-3$ for zero-defect matchings, and it isolates the terms $k=4$ and $k=n-4$; together these turn the source's $O(n^{-1})$ into $O(n^{-2})$. The count \eqref{cgr:fc:eq:source-decomp} is its Lemma~5.6 (both checked in arXiv:2402.06394v1). \eqref{cgr:fc:eq:prime-prob} gives Part~I's \eqref{cgr:ld:eq:prime-prob} its first correction.
\end{writenote}

\section{A graph level finite split tree formula}\label{cgr:fc:sec:split}
Let $r_s$ be the number of connected circle-graph isomorphism
classes with one distinguished vertex. Rooting means choosing a
vertex orbit, so \eqref{cgr:fc:eq:coarse} gives
\begin{equation}\label{cgr:fc:eq:root-bound}
 r_s\le s c_s\le C m_s.
\end{equation}
Put $b_j=r_{j+2}$ and $B_d(z)=\sum_{j=0}^d b_jz^j$.

\begin{theorem}[Finite canonical transfer]\label{cgr:fc:thm:finite-transfer}
For every fixed integer $d\ge1$,
\begin{equation}\label{cgr:fc:eq:finite-transfer}
 c_n=\sum_{N=\ceil{n/(d+1)}}^n p_N^{\mathrm a}
           [z^{n-N}]B_d(z)^N+O(h_n n^{-d}).
\end{equation}
Apart from a superpolynomially negligible symmetric-core class,
the sum counts exactly the connected graphs with no split whose
two sides both have at least $d+2$ vertices, for all sufficiently
large $n$.
\end{theorem}

\subsection{Large split errors without choosing a representation}
A split of a connected circle graph with sides $s,n-s$ produces rooted factors of sizes
$s+1,n-s+1$ by adjoining a marker adjacent to the frontier on
its side. Both factors are connected. Each is isomorphic to an
induced subgraph of the original circle graph: use any vertex in
the opposite nonempty frontier for the marker. The two rooted
isomorphism classes determine their split composition uniquely,
by deleting both roots and joining their neighbor sets completely.
Uniqueness of the original split is unnecessary for an upper bound.

Let $E_{n,r}$ count connected circle graphs with a split whose sides both have at
least $r\ge2$ vertices. Then
\begin{equation}\label{cgr:fc:eq:large-split}
 E_{n,r}\le\sum_{s=r}^{\floor{n/2}}r_{s+1}r_{n-s+1}
 \le C^2\sum_{s=r}^{\floor{n/2}}m_{s+1}m_{n-s+1}
 =O(h_n n^{-r+2}).
\end{equation}
For the last equality put $a_s=m_{s+1}m_{n-s+1}$.
The ratio $a_{s+1}/a_s=(2s+3)/(2n-2s+1)$ is at most one in
the needed half range, so the sum is bounded by
$a_r+na_{r+1}=O(a_r)=O(m_{n-r+1})$ for fixed $r$.
Finally $m_{n-r+1}/m_n=O(n^{-r+1})$.
In particular, every graph with a split whose sides both have at
least four vertices contributes in total $O(h_n/n^2)$.
The proof operates on graph factors, not interval representations.

\subsection{The canonical large prime center}
Fix $r\ge3$ and suppose no graph split has two sides of size at
least $r$. Give each leaf of the canonical reduced split tree
weight one and each internal node weight zero. A weighted centroid
exists with every component after its removal containing at most
$n/2$ leaves: repeatedly move towards a component of weight
larger than $n/2$. This process advances along the finite tree,
so stops, and for $n>2$ stops at an internal node $v$.
Every component at $v$ has at most $r-1$ leaves. Otherwise its
internal attachment edge induces a forbidden split; a component
attached by a leaf edge has size one. If $v$ has degree $N$, then
\begin{equation}\label{cgr:fc:eq:core-size}
 N\ge n/(r-1).
\end{equation}
For $n\ge3r-2$, $v$ cannot be degenerate. Greedily group its
ports until their components total between $r$ and $2r-2$
leaves. This group and its complement each use at least two ports,
and each carry at least $r$ leaves. Any partition of a clique or
star label into two sets of at least two markers is a split.
Splitting that degenerate node would thus give a forbidden graph
split, a contradiction.

For $n>r(r-1)$ the prime center has degree $N>r$.
Every other internal node is inside a branch of at most $r-1$
leaves. Adjoining its attachment marker gives at most $r$ leaves,
so such a node has degree at most $r$, since each of its incident
components contains a leaf. The center is therefore the unique
node of degree greater than $r$ and is canonically recovered by
graph isomorphisms.

Conversely, attach branches of at most $r-1$ actual vertices to
a prime center of degree $N>r$. By source Corollary 2.10, any
split comes from a tree edge or one degenerate-node split. The
prime center cannot split; every degenerate node is inside one
branch. In each case the side away from the center has at most
$r-1$ leaves. Thus no split has two sides of size at least $r$.
This proves the needed characterization, including its converse.

\subsection{Rooted branches and isomorphism multiplicities}
For each port $x$ of a prime graph $H$, choose a rooted connected
circle graph $(S_x,\rho_x)$ of order $j+2$ with $0\le j\le d$.
Its actual branch vertices are $V(S_x)\setminus\{\rho_x\}$,
and its excess over the core port is $j$. With
$A_x=N_{S_x}(\rho_x)$, retain all within-branch edges and put all
edges of $A_x\times A_y$ between branches exactly when $xy$ is
an edge of $H$. This is simultaneous split composition.

The identity branch is rooted $K_2$. There is no order-one branch,
which would leave zero actual vertices after root deletion, and no
second distinguished attachment vertex or marker orientation.
At the tree level, remove the root leaf from the reduced tree of
$S_x$ and connect its former neighbor to the center port. For
$K_2$ simply attach the actual leaf. All old adjacencies remain
reduced, and each new internal adjacency has a prime endpoint;
therefore no reduction changes the core or the branch classes.
Cutting a branch and restoring its root leaf recovers exactly
$(S_x,\rho_x)$.

If $H$ is asymmetric, assignments give the same unmarked graph
exactly when each rooted branch type agrees at the same port.
Thus the exact contribution of $H$ is $[z^{n-N}]B_d(z)^N$.
There is no division by $N!$, no extra root count, and no division
by branch automorphism groups: rooted isomorphism types already
account for them. Indeed
\[
 \Aut(G)\cong\prod_{x\in V(H)}\Aut(S_x,\rho_x)
 \quad\hbox{when }\Aut(H)=1.
\]
For a symmetric prime center the exact count would require its
port-action orbits; the upper bound $B_d(1)^N$ suffices here.
Using \eqref{cgr:fc:eq:symmetry}, all such contributions are at most
\begin{equation}\label{cgr:fc:eq:symcore}
 \sum_{N\le n}S_N B_d(1)^N
 \le\exp\bigl(\tfrac12 n\log n+O_d(n)\bigr)
 =o(h_n n^{-A})
\end{equation}
for every fixed $A$. Combining the canonical characterization,
\eqref{cgr:fc:eq:large-split} with $r=d+2$, and \eqref{cgr:fc:eq:symcore}
proves Theorem~\ref{cgr:fc:thm:finite-transfer}.

\subsection{The first rooted counts}\label{cgr:fc:sub:rooted}
Only rooted $K_2$ occurs at order two, so $b_0=1$.
At order three, $P_3$ has two root orbits and $K_3$ has one,
giving $b_1=3$. These are the leaf, false-twin and true-twin
one-vertex decorations. The six connected graphs of order four
have the following root-orbit counts:
\[
 P_4:2,\quad K_{1,3}:2,\quad C_4:1,\quad
 \text{paw}:3,\quad\text{diamond}:2,\quad K_4:1.
\]
Each is a circle graph, so $b_2=11$.
The companion independently enumerates their graph classes,
representations and root orbits. These finite facts determine the
first correction; their finiteness does not supply any analytic
uniformity statement.
\begin{writenote}[Checked, 5~October 2026]
A brute-force enumeration by the write over all graphs on two to five vertices gives $1$, $2$, $6$, $21$ connected classes and $1$, $3$, $11$, $58$ vertex orbits, and root-stabilizer orders $\{1{:}1\}$, $\{1{:}1,2{:}2\}$, $\{1{:}3,2{:}6,6{:}2\}$ at excess $0$, $1$, $2$, as used in Subsection~\ref{cgr:fc:sub:groups}. That every graph on at most five vertices is a circle graph is the companion's enumeration; the write did not repeat it.
\end{writenote}

\section{Global tails and the first graph coefficient}\label{cgr:fc:sec:coeff}
\subsection{Uniform factorial domination}
A branch multiplicity vector $(a_1,\ldots,a_d)$ has excess
$k=\sum_j ja_j$, occupied-port count $\ell=\sum_ja_j$, and core
order $N=n-k$. Its contribution to \eqref{cgr:fc:eq:finite-transfer} is
\begin{equation}\label{cgr:fc:eq:branch-term}
 p_N^{\mathrm a}(N)_\ell\prod_{j=1}^d\frac{b_j^{a_j}}{a_j!},
 \qquad \ell\le N,
\end{equation}
where $(N)_\ell$ is a falling factorial. Since $n\le(d+1)N$,
\begin{align}
 \frac{h_N}{h_n}(N)_\ell
 &=\frac nN\frac{m_N}{m_n}(N)_\ell\notag\\
 &\le(d+1)(2N)^{-k}N^\ell\notag\\
 &\le(d+1)\prod_{j=1}^d
 \left[2^{-j}\left(\frac{d+1}{n}\right)^{j-1}\right]^{a_j}.
 \label{cgr:fc:eq:factorial-dom}
\end{align}
Using $p_N^{\mathrm a}=O(h_N)$, the normalized term is bounded by
a constant times $\prod_j\theta_j(n)^{a_j}/a_j!$, where
\begin{equation}\label{cgr:fc:eq:theta}
 \theta_j(n)=b_j2^{-j}\bigl((d+1)/n\bigr)^{j-1},
 \qquad \theta_1=3/2.
\end{equation}
Vectors with $\sum_{j\ge2}(j-1)a_j\ge q+1$ have total weight
$O(n^{-q-1})$: factor out that inverse power and sum the fixed
exponential weights. The factorial tail $a_1>\ceil{\log n}$
is smaller than every fixed inverse power. Both bounds apply
across the full range of core orders down to $n/(d+1)$.
A fixed-deficit expansion is used only after this tail control.

\subsection{First coefficient for connected graphs}
Suppose temporarily
\begin{equation}\label{cgr:fc:eq:prime-alpha}
 p_n^{\mathrm a}=P h_n(1+\alpha/n+O(n^{-2})),\quad P>0.
\end{equation}
Take $d=2$. By \eqref{cgr:fc:eq:theta}, terms with two or more
excess-two branches are $O(h_n/n^2)$. Write $a=a_1,b=a_2$,
and restrict $a\le\ceil{\log n}$ at a superpolynomial error.
For $b=0$, exact cancellation gives
\begin{align}
 \frac{3^a(n-a)_a h_{n-a}}{a!h_n}
 &=\frac{(3/2)^a}{a!}\frac n{n-a}
   \prod_{j=0}^{a-1}
   \frac{1-(a+j)/n}{1-(j+1/2)/n}\notag\\
 &=\frac{(3/2)^a}{a!}
 \left(1+\frac{-a^2+3a/2}{n}
       +O\left(\frac{(a+1)^4}{n^2}\right)\right).
 \label{cgr:fc:eq:onebranch-exp}
\end{align}
Taylor's formula for $\log(1-x)$ is uniform on this logarithmic
range; sums of squared offsets are $O((a+1)^3)$, and the squared
linear coefficient is $O((a+1)^4)$.
The prime multiplier is
$P(1+\alpha/n+O((a+1)/n^2))$. For $A\sim\Poi(3/2)$,
$\EE[-A^2+3A/2]=-3/2$. The sum of the $b=0$ terms is thus
\[
 P e^{3/2}h_n\left(1+\frac{\alpha-3/2}{n}+O(n^{-2})\right).
\]
For $b=1$, $N=n-a-2$, $\ell=a+1$, and the analogous ratio to
leading order is
\[
 \frac{p_N^{\mathrm a}3^a11(N)_{a+1}}{a!h_n}
 =P\frac{(3/2)^a}{a!}\frac{11}{4n}
   \left(1+O((a+1)^2/n)\right).
\]
Summing proves
\begin{equation}\label{cgr:fc:eq:connected-transfer}
 c_n=P e^{3/2}h_n
       \left(1+\frac{\alpha+5/4}{n}+O(n^{-2})\right).
\end{equation}
With $P=e^{-3}$ and $\alpha=-5$ from \eqref{cgr:fc:eq:prime-main},
this is \eqref{cgr:fc:eq:c-main}. The elementary Stirling expansion
$m_n=\sqrt2(2n/e)^n(1-1/(24n)+O(n^{-2}))$ gives its alternate
normalization in Theorem~\ref{cgr:fc:thm:main}.

\section{Components and every fixed order}\label{cgr:fc:sec:allorders}
\subsection{The component transfer with a proved tail}
For every fixed $q\ge0$,
\begin{equation}\label{cgr:fc:eq:component-transfer}
 g_n=\sum_{k=0}^qg_kc_{n-k}+O(h_n n^{-q-1}).
\end{equation}
For $n>2q$, the displayed terms count exactly graphs with a
component of order at least $n-q$, uniquely. Any other graph can
have its whole connected components grouped into two sets of
orders $s,n-s$ with $q+1\le s\le n/2$, for sufficiently large
$n$. Take a component already in that interval, the complement
of a larger component, or greedily combine components all of
order at most $q$. For $q=0$ a disconnected graph directly gives
such a grouping. Overcounting yields the upper bound
\[
 \sum_{s=q+1}^{\floor{n/2}}g_sg_{n-s}
 \le\frac{2C^2}{n}\sum_{s=q+1}^{\floor{n/2}}m_sm_{n-s}
 =O(m_{n-q-1}/n)=O(h_n n^{-q-1}).
\]
The last step is the same endpoint-dominant ratio argument as in
\eqref{cgr:fc:eq:large-split}. No automorphism factor is needed for a
unique large component and its disjoint remaining graph.

At $q=1$, $g_1=1$ gives
$g_n=c_n+c_{n-1}+O(h_n/n^2)$, and
$h_{n-1}/h_n=1/(2n)+O(n^{-2})$. Thus
\begin{equation}\label{cgr:fc:eq:all-transfer}
 g_n=P e^{3/2}h_n
       \left(1+\frac{\alpha+7/4}{n}+O(n^{-2})\right),
\end{equation}
proving \eqref{cgr:fc:eq:g-main}.

\subsection{A second connectivity term without a second prime term}
Write $A=P e^{3/2}$, $\gamma=\alpha+5/4$ and
$\eta=\gamma+1/2$. At $q=2$, since $g_2=2$,
\[
 g_n-c_n=c_{n-1}+2c_{n-2}+O(h_n/n^3).
\]
The factorial ratios are
\[
 \frac{h_{n-1}}{h_n}=\frac1{2n}+\frac3{4n^2}+O(n^{-3}),
 \qquad \frac{h_{n-2}}{h_n}=\frac1{4n^2}+O(n^{-3}).
\]
Using only \eqref{cgr:fc:eq:connected-transfer},
\[
 \frac{g_n-c_n}{A h_n}
 =\frac1{2n}+\frac{5/4+\gamma/2}{n^2}+O(n^{-3}).
\]
Divide by $g_n/(A h_n)=1+\eta/n+O(n^{-2})$ to obtain
\[
 \frac{g_n-c_n}{g_n}
 =\frac1{2n}+\frac{5/4+\gamma/2-\eta/2}{n^2}+O(n^{-3})
 =\frac1{2n}+\frac1{n^2}+O(n^{-3}).
\]
The extra $1/n$ in the smaller-component ratio makes first-order
individual expansions sufficient for the $O(n^{-3})$ remainder.
No existence of an individual second coefficient is assumed.

\subsection{The conditional transfer theorem}\label{cgr:fc:sub:conditional}
Suppose for some fixed $q$ that
\begin{equation}\label{cgr:fc:eq:prime-q}
 \frac{p_n^{\mathrm a}}{h_n}
 =P\sum_{t=0}^q\alpha_t n^{-t}+O(n^{-q-1}),\qquad\alpha_0=1.
\end{equation}
Use Theorem~\ref{cgr:fc:thm:finite-transfer} with $d=q+1$.
The graph-theoretic error is $O(h_n n^{-q-1})$.
Only vectors with $\sum_{j\ge2}(j-1)a_j\le q$ matter; there
are finitely many such choices of the coordinates $j\ge2$.
For each, restrict $a_1\le\ceil{\log n}$ by the global bound.
The exact factorial ratio in \eqref{cgr:fc:eq:branch-term} has a Taylor
expansion whose coefficients are polynomials in $a_1$ and whose
remainder is $n^{-q-1}$ times a fixed polynomial, multiplied by
$(3/2)^{a_1}/a_1!$. The same is true after multiplying by
\eqref{cgr:fc:eq:prime-q}; replacing $n-k$ by $n$ is uniform on this
range. All polynomial Poisson moments are finite, so termwise
summation is justified. This proves the connected transfer through
order $q$, and \eqref{cgr:fc:eq:component-transfer} proves the all-graph
transfer. Only $b_0,\ldots,b_{q+1}$ and $g_0,\ldots,g_q$ enter.
The statement is for each fixed order, with no convergence claim
and no error uniform as $q$ grows.

\subsection{Conditional second coefficients}\label{cgr:fc:sub:second}
For clarity, the next transfer calculation can be made explicit.
Assume the additional prime hypothesis
\[
 p_n^{\mathrm a}/h_n=P(1+\alpha/n+\beta/n^2+O(n^{-3})).
\]
All 21 connected graph classes on five vertices are circle graphs;
their root-orbit count is $b_3=58$, checked by complete finite
enumeration in the companion. For excess $k$ and port count $\ell$,
\begin{align*}
 \frac{h_{n-k}}{h_n}(n-k)_\ell
 &=2^{-k}n^{\ell-k}
 \left(1+\frac{L_1}{n}+\frac{L_2+L_1^2/2}{n^2}
            +O(n^{-3}\operatorname{poly}(k+\ell))\right),\\
 L_1&=k+k^2/2-\ell k-\ell(\ell-1)/2,\\
 L_2&=k^2/2+k(4k^2-1)/24\\
 &\hspace{7mm}-\bigl(\ell k^2+k\ell(\ell-1)
            +\ell(\ell-1)(2\ell-1)/6\bigr)/2.
\end{align*}
Here the remainder has the logarithmic-range meaning just proved,
not a claim uniform in unrestricted $k,\ell$.
For only $a$ excess-one branches, the first two polynomials are
\[
 A_1=-a^2+3a/2,\qquad
 A_2=a^4/2-5a^3/2+19a^2/8-a/8.
\]
For one excess-two branch they require only the first polynomial
$B_1=-a^2-a/2+2$.
For $A\sim\Poi(3/2)$,
\[
 \EE A_1=-3/2,\quad \EE A_2=15/16,\quad
 \EE(A_1+A)=0,\quad \EE B_1=-5/2.
\]
The other retained possibilities are two excess-two branches,
with leading weight $121/(32n^2)$, and one excess-three branch,
with weight $58/(8n^2)$. Including the shifted prime multiplier
$1+\alpha/n+(\alpha k+\beta)/n^2$ gives
\begin{align}
 \frac{c_n}{P e^{3/2}h_n}
 &=1+\frac{\alpha+5/4}{n}
   +\frac{\beta+11\alpha/4+163/32}{n^2}+O(n^{-3}),
 \label{cgr:fc:eq:secondc}\\
 \frac{g_n}{P e^{3/2}h_n}
 &=1+\frac{\alpha+7/4}{n}
   +\frac{\beta+13\alpha/4+223/32}{n^2}+O(n^{-3}).
 \label{cgr:fc:eq:secondg}
\end{align}
The second line adds $\gamma/2+5/4$ by the component identity.
With $\alpha=-5$ these yield the coefficients stated in
Theorem~\ref{cgr:fc:thm:transfer-summary}, still conditional on $\beta$.
\begin{writenote}[Checked numerically, 5~October 2026]
The write evaluated the model sum $\sum_NPh_N(1+\alpha/N+\beta/N^2)[z^{n-N}]B_3(z)^N$, with $b_0,\ldots,b_3=1,3,11,58$ and $n-N\le120$, at $n=1000$, $2000$, $4000$, $8000$ in 40-digit arithmetic. Richardson extrapolation of $n^2$ times its second-order residual reproduces $\beta+\frac{11\alpha}4+\frac{163}{32}$ to within $10^{-5}$ for $(\alpha,\beta)=(0,0)$, $(-5,0)$, $(1,0)$, $(0,1)$; in particular $-\frac{277}{32}=-8.65625$ at $\alpha=-5$, $\beta=0$. By hand, $g_n=c_n+c_{n-1}+2c_{n-2}+O(h_n/n^3)$ with $h_{n-1}/h_n=\frac1{2n}+\frac3{4n^2}+O(n^{-3})$ and $h_{n-2}/h_n=\frac1{4n^2}+O(n^{-3})$ turns \eqref{cgr:fc:eq:secondc} into \eqref{cgr:fc:eq:secondg}. This checks the algebra of the two displays from the formula of Theorem~\ref{cgr:fc:thm:finite-transfer}; it does not check the transfer theorem, and it says nothing about~$\beta$.
\end{writenote}

\section{Bounded automorphism weights}\label{cgr:fc:sec:weighted}
The same canonical proof yields additional first corrections for
bounded observables. For $t\in[0,\infty)$, define
\[
 c_n(t)=\sum_{\substack{G\text{ connected}\\|G|=n}}|\Aut(G)|^{-t},
 \qquad g_n(t)=\sum_{|G|=n}|\Aut(G)|^{-t}.
\]
Both sums range over unlabeled circle-graph isomorphism classes.
At $t=\infty$, interpret the weight as the indicator of a trivial
automorphism group. Every weight lies in $[0,1]$, so excluded
large-split and symmetric-core classes still have mass
$O(h_n/n^2)$, uniformly in $t$. On the recovered asymmetric-core
class, the automorphism group is the direct product of the rooted
branch groups. Replace each $b_j$ by
\[
 b_j(t)=\sum_{\substack{(S,\rho)\text{ rooted connected circle graph}\\
                       |S|=j+2}}|\Aut(S,\rho)|^{-t}.
\]
Then $0\le b_j(t)\le b_j(0)$ and all the factorial bounds remain
uniform. This is a bounded-weight argument; it does not control
positive powers of the automorphism-group order.

\subsection{Weighted transfer and small branch groups}\label{cgr:fc:sub:groups}
With general $b_1,b_2$ the Poisson parameter is $b_1/2$.
The expectation of $-A^2+3A/2$ is
$-b_1^2/4+b_1/4$, and one excess-two branch contributes $b_2/4$.
Consequently, uniformly in $t\in[0,\infty]$,
\begin{equation}\label{cgr:fc:eq:weightedc}
 c_n(t)=e^{-3+b_1(t)/2}h_n
 \left(1+\frac{-5+(b_2(t)+b_1(t)-b_1(t)^2)/4}{n}
              +O(n^{-2})\right).
\end{equation}
The uniformity follows from the already uniform graph error, the
prime error independent of $t$, and bounded Poisson moments for
$0\le b_1(t)\le3$. In an all-graph decomposition the unique large
component cannot be exchanged with its smaller complement, so the
automorphism weight multiplies. The singleton complement has
weight one, giving $g_n(t)=c_n(t)+c_{n-1}(t)+O(h_n/n^2)$;
its relative coefficient in \eqref{cgr:fc:eq:weightedc} increases by $1/2$.

The three excess-one rooted branch groups have orders $1,2,2$.
The eleven excess-two groups have orders $1$ in three cases,
$2$ in six cases, and $6$ in two cases. The rigid cases are the
two root orbits of $P_4$ and the paw rooted at an interchangeable
degree-two vertex. The order-six cases are the star rooted at its
center and the clique $K_4$, both with group $S_3$; all order-two
groups are $C_2$. Thus
\begin{equation}\label{cgr:fc:eq:weighted-b}
 b_1(t)=1+2\cdot2^{-t},\qquad
 b_2(t)=3+6\cdot2^{-t}+2\cdot6^{-t}.
\end{equation}
These small group counts are exhaustively checked in the companion.

\subsection{Labeled graphs and asymmetry}
At $t=1$, $b_1=2,b_2=19/3$, and
$(b_2+b_1-b_1^2)/4=13/12$. Multiplying the reciprocal-
automorphism sum by $n!$ counts labeled graphs on a prescribed
$n$-element vertex set. Therefore
\begin{align}
 C_n^{\mathrm{lab}}&=n!e^{-2}h_n
          \left(1-\frac{47}{12n}+O(n^{-2})\right),\label{cgr:fc:eq:labeledc}\\
 G_n^{\mathrm{lab}}&=n!e^{-2}h_n
          \left(1-\frac{41}{12n}+O(n^{-2})\right).\label{cgr:fc:eq:labeledg}
\end{align}
These are mathematical consequences of the bounded-weight proof;
no separate labeled OEIS identification or exhaustive labeled
literature-priority audit is claimed.

At $t=\infty$, $b_1=1,b_2=3$, so
\[
 c_n(\infty)=e^{-5/2}h_n\left(1-\frac{17}{4n}+O(n^{-2})\right),
 \quad
 g_n(\infty)=e^{-5/2}h_n\left(1-\frac{15}{4n}+O(n^{-2})\right).
\]
Under either uniform connected or uniform arbitrary unlabeled
sampling, division by the respective ordinary count gives
\begin{equation}\label{cgr:fc:eq:asymmetric}
 \PP(\Aut(G)=1)=e^{-1}\left(1-\frac1{2n}+O(n^{-2})\right).
\end{equation}
More generally put $\rho=2^{-t}$ and $\sigma=6^{-t}$, with both
zero at infinity. Equations \eqref{cgr:fc:eq:weightedc}--\eqref{cgr:fc:eq:weighted-b}
give, uniformly in $t\in[0,\infty]$ under either sampling law,
\begin{equation}\label{cgr:fc:eq:aut-transform}
 \EE[|\Aut(G)|^{-t}]
 =e^{\rho-1}\left(1+
       \frac{-1/2+\rho-\rho^2+\sigma/2}{n}+O(n^{-2})\right).
\end{equation}
At $t=1$ this is
$e^{-1/2}(1-1/(6n)+O(n^{-2}))$.
The universal prime coefficient cancels from the normalized
transform, although its controlled first-order expansion underlies
the argument.

\subsection{Rare nonabelian groups}
Through first order, all recovered branch groups are trivial or
$C_2$, apart from the two excess-two branches with group $S_3$.
Two excess-two branches, larger branch types and excluded splits
have probability $O(n^{-2})$. Replacing the single excess-two
count $11$ by $2$ in Section~\ref{cgr:fc:sec:coeff} gives
\begin{equation}\label{cgr:fc:eq:nonabelian}
 \PP(\Aut(G)\text{ nonabelian})=\frac1{2n}+O(n^{-2}).
\end{equation}
The same estimate holds for not being an elementary abelian
$2$-group. For each fixed $k\ge0$, the refined law is
\begin{equation}\label{cgr:fc:eq:s3-law}
 \PP\bigl(\Aut(G)\cong S_3\times(C_2)^k\bigr)
 =\frac{e^{-1}}{2n\,k!}+O(n^{-2})
\end{equation}
under either connected or arbitrary unlabeled sampling.
Indeed, the two $C_2$ excess-one branch types have total limiting
Poisson mean one, the rigid type has mean $1/2$, and one $S_3$
branch contributes $2/(4n)$. Summing rigid branches and selecting
exactly $k$ nonrigid excess-one branches gives the factor
$e^{-1}/k!$ after normalization. For arbitrary sampling, the only disconnected class at order
$1/n$ is a giant connected graph plus one isolated vertex. Its
nonabelian subfamily pays a further $1/n$ factor and so has
probability $O(n^{-2})$; all other disconnected exceptions are
already $O(n^{-2})$. This justifies the same rare-group coefficients
for arbitrary sampling. This fixed-$k$ statement is not an
unbounded automorphism moment claim.

\section{Refined least index inverses}\label{cgr:fc:sec:inverse}
For $a_n=c_n$ or $g_n$, define
\[
 N_a(y)=\min\{n\ge1:a_n\ge y\}.
\]
Both sequences are unbounded and eventually strictly increasing.
Indeed Theorem~\ref{cgr:fc:thm:main} gives
$a_{n+1}/a_n\sim h_{n+1}/h_n\sim2n$, without differentiating any
unknown error. Their finite pre-onset prefixes are bounded, so for
sufficiently large $y$ the least threshold lies in the increasing
tail. This is enough for the asymptotic inverse; no separate
finite-order monotonicity assertion is required.

Let $L=\log y$, and use the positive real Lambert branch $W_0$:
\begin{equation}\label{cgr:fc:eq:inverse-vars}
 u=\frac{L}{W_0(2L/e)},\qquad D=\log(2u),\qquad
 B=\frac{3+\log2}{2}.
\end{equation}
Then $u(\log(2u)-1)=L$. For the connected and all-graph cases
respectively set
\[
 K_c=\frac{103}{24},\qquad K_g=\frac{91}{24},\qquad
 s_a(y)=u+1+\frac BD+
       \frac{K_a-B^2/(2D^2)}{uD}.
\]

\begin{theorem}[Ceiling safe first refinement]\label{cgr:fc:thm:inverse}
For each of $a=c,g$, there exist constants $C_a,y_a>0$ such that,
for every $y\ge y_a$,
\begin{equation}\label{cgr:fc:eq:inverse-bracket}
 \ceil{s_a(y)-\frac{C_a}{u^2D}}
 \ \le\ N_a(y)\ \le\
 \ceil{s_a(y)+\frac{C_a}{u^2D}}.
\end{equation}
The centers satisfy $s_c(y)-s_g(y)=1/(2uD)$.
The bracket does not authorize unconditional rounding at inputs
where its interval meets an integer. Its constants and onset are
noneffective in this proof.
\end{theorem}

\begin{proof}
Write $a_n=e^{-3/2}h_n(1+\delta/n+O(n^{-2}))$, with
$\delta=-15/4$ or $-13/4$. The smooth positive normalization is
\[
 H(t)=\frac{e^{-3/2}2^t\Gamma(t+1/2)}{4\sqrt\pi\,t},
 \qquad F_\delta(t)=\log H(t)+\frac\delta t.
\]
Stirling's expansion gives
\begin{equation}\label{cgr:fc:eq:smoothlog}
 \log H(t)=t(\log(2t)-1)-\log t
 -\frac{3(1+\log2)}2-\frac1{24t}+O(t^{-3}).
\end{equation}
Its derivative is obtained from the smooth gamma function, so
$F_\delta'(t)=\log(2t)-1/t+O(t^{-2})>0$ eventually.
Also $\log a_n=F_\delta(n)+O(n^{-2})$; no derivative is taken
of this integer-sequence remainder.

For bounded $d$, expanding at $u$ gives
\begin{align*}
 F_\delta(u+d)-L
 =Dd-\log u-\frac{3(1+\log2)}2
 +\frac{d^2/2-d-1/24+\delta}{u}+O(u^{-2}).
\end{align*}
Since $\log u=D-\log2$, the order-one solution is
$d_0=1+B/D$. Put $d=d_0+d_1/u$. The next coefficient vanishes
when
\[
 d_1=\frac{13/24-\delta-B^2/(2D^2)}D.
\]
Thus $K_a=13/24-\delta$, giving the claimed values. Substitution
has residual $O(u^{-2})$. The unique large root of
$F_\delta(t)=L$ differs from $s_a(y)$ by $O(1/(u^2D))$,
because the derivative is comparable with $D$ near $u$.

To pass to the integer sequence, first a coarse leading bracket
locates the threshold and the smooth roots in $[u/2,2u]$ for
large $y$. The tail supremum of
$|\log a_n-F_\delta(n)|$ over integers $n\ge\floor{u/2}$ is
$O(u^{-2})$. Let $r_-$ and $r_+$ solve
$F_\delta(r_\pm)=L\pm C'u^{-2}$, with $C'$ large enough.
If $n<\ceil{r_-}$ in this tail then $F_\delta(n)<L-C'u^{-2}$,
so $a_n<y$; at $n=\ceil{r_+}$, $a_n\ge y$.
The bounded prefix is already below $y$. Hence
$\ceil{r_-}\le N_a(y)\le\ceil{r_+}$.
Both roots differ from $s_a(y)$ by $O(1/(u^2D))$, proving
\eqref{cgr:fc:eq:inverse-bracket} after increasing its constant.
\end{proof}

Higher prime coefficients, if proved with controlled errors, could
be passed through the fixed-order transfer and then inverted
recursively by the same smooth procedure. The displayed theorem
uses only the unconditional first corrections. It makes no
unconditional all-orders inverse claim.
\begin{writenote}[The inverse, 5~October 2026]
At $x=y$, $s_a(y)=s(x)+(K_a-B^2/(2D^2))/(uD)$, with $s$ Part~I's centre \eqref{cgr:ld:eq:s}. The leading $u$ is an instance of the transseries volume's factorial core, and the bracket follows the pattern of its staircase theorem without being an instance (front matter).
\end{writenote}

\section{Reproducibility and limits}\label{cgr:fc:sec:checks}
The companion is a fresh exact implementation for this report.
Its role is to verify finite atom, graph and rational identities,
not to certify analytic asymptotics from samples. In particular it
checks local overlap and compatible-cluster counts, the four
three-chord interval cases, rooted graph-isomorphism counts,
first and conditional second transfer arithmetic, the inverse
coefficient cancellation, and small rooted automorphism orders. The README records the exact tested
ranges and outputs. Explicit exceptions, rather than removable
Python assertions, enforce the checks. Mutation tests require
incorrect coefficients to fail, including under Python's
optimization mode; output-safety tests reject overwrites and
symlink routes.

The package supplies editable LaTeX, the rendered PDF, exact
companion source and evidence, source provenance, checksums,
deterministic builders and release tests. It excludes downloaded
papers, exploratory working notes, and earlier report files.
PDF byte reproducibility is qualified to the tested pdfTeX
installation; the mathematics does not depend on identical PDF
bytes. ZIP byte reproducibility uses fixed entry order, timestamps,
permissions and stored payloads. A checksum manifest checks
consistency, not authenticity against joint replacement of a file
and its hash.

The established new content is the uniform local first correction,
its direct $O(n^{-2})$ nonlocal comparison, the canonical finite
split-tree transfer, the individual graph first coefficients, the
second connectivity term, and the refined ceiling-safe inverse.
An unknown prime second coefficient remains unknown. Effective
constants, a final-journal source reconciliation, and any stronger
priority claim remain outside the proved scope. The bounded reciprocal-automorphism refinements and labeled
consequences in Section~\ref{cgr:fc:sec:weighted} are included with their
proof; unbounded automorphism moments are not claimed.
\begin{writenote}[Status, 5~October 2026]
The unknown $\beta$ and the other limits stated here are collected, with Part~I's, in Section~\ref{cgr:sec:further}. The companion is shipped as \file{code/152-corr-companion-*.py} and \file{data/152-corr-companion-*}; at intake its checks were rerun on a copy and the recorded result regenerated byte for byte (README). The \file{SOURCE_PROVENANCE.json} entries on ``independent review'' are metadata, as that file says, not an independent proof certificate.
\end{writenote}

\section{Further questions and research}\label{cgr:sec:further}
\begin{writenote}[Standing rule, 5~October 2026]
This section is added in the write. Under Vladimir's standing rule of 4~October 2026 it collects every question left open by the two manuscripts, with its source, the argument sketch it came with, and what is missing. Part~I's Section~\ref{cgr:ld:sec:limits} item~1 is answered by Part~II (dated note there). Nothing in the two manuscripts was found to be wrong, and no claim was refuted.
\end{writenote}
\begin{enumerate}
\item\label{cgr:q:beta} \emph{The prime second coefficient $\beta$} (Part~II, Theorem~\ref{cgr:fc:thm:transfer-summary}, Subsection~\ref{cgr:fc:sub:second}, Section~\ref{cgr:fc:sec:inverse}). If $p_n^{\mathrm a}=e^{-3}h_n(1-\frac5n+\frac\beta{n^2}+O(n^{-3}))$, then the second relative coefficients of $c_n$ and $g_n$ are $\beta-\frac{277}{32}$ and $\beta-\frac{297}{32}$ (proved conditionally), and the $n^{-2}$ term of $c_n/g_n$ does not depend on $\beta$ (proved). Part~II gives no value and no conjecture. What is missing is the next order in each step of Sections~\ref{cgr:fc:sec:polymer}--\ref{cgr:fc:sec:matching}: the polymers of four chords, whose total activity Subsection~\ref{cgr:fc:sec:polymer-bounds} bounds by $O(M^{-2})$, the hard-core expansion \eqref{cgr:fc:eq:hardcore} beyond its pair term, the matching probabilities \eqref{cgr:fc:eq:q-global} beyond $r^2/M$, and the nonlocal term: zero-defect matchings with a $k$-decomposition, $4\le k\le n-4$, have probability $O(n^{-2})$ by the bound before \eqref{cgr:fc:eq:prime-prob}, so at second order they must be counted rather than bounded. The intake Monte Carlo concerns only the first-order constant and cannot resolve a second-order term at its sample sizes.
\item\label{cgr:q:allorders} \emph{An unconditional expansion to every fixed order} (Part~I, Section~\ref{cgr:ld:sec:limits} item~2; Part~II, Theorem~\ref{cgr:fc:thm:transfer-summary} and Subsection~\ref{cgr:fc:sub:conditional}). Part~II reduces every fixed order of $c_n$ and $g_n$ to the same order of $p_n^{\mathrm a}/h_n$. It sketches that higher prime coefficients ``could be passed through the fixed-order transfer and then inverted recursively by the same smooth procedure'' (Section~\ref{cgr:fc:sec:inverse}); what is missing is the prime expansion itself at each order (item~\ref{cgr:q:beta} is its first case). Neither Part claims convergence, a remainder uniform in the order, or an all-orders inverse.
\item\label{cgr:q:effective} \emph{Effective constants, onsets and certified thresholds} (Part~I, item~3 and the remarks after Theorem~\ref{cgr:ld:thm:quantitative}; Part~II, scope paragraph and Theorem~\ref{cgr:fc:thm:inverse}). Only the local bound $\min\{1,200/n\}$ of \eqref{cgr:ld:eq:explicit-TV} is explicit. The graph-count constants, the constant of \eqref{cgr:ld:eq:uniform-g} and \eqref{cgr:fc:eq:coarse}, the onsets, and the inverse constants $C_*$, $C_a$ exist but are not computed, so neither bracket decides the least index at a given input. What is missing is an explicit constant in every $O$-term of one complete chain; Part~II's chain does not use the qualitative limit \eqref{cgr:ld:eq:poi-source}, so it is the natural candidate.
\item\label{cgr:q:sources} \emph{Source reconciliation and priority} (Part~I, item~4 and Subsection~\ref{cgr:ld:sub:inspected}; Part~II, Subsection~\ref{cgr:fc:sub:priority}). The 2026 journal text of \cite{BBFGP24} was not read at intake or at the write; Gabor, Supowit and Hsu's paper was read only through \cite[Proposition 5.12]{BBFGP24}, and the paper of Arratia, Bollob\'as, Coppersmith and Sorkin at abstract level. A stronger priority statement than ``not located in the sources checked'' needs these.
\item\label{cgr:q:labeled} \emph{Labeled counts} (Part~I, item~4 and the paragraph after \eqref{cgr:ld:eq:labeled}; Part~II, after \eqref{cgr:fc:eq:labeledg}). The labeled equivalents and first corrections are proved; a literature and OEIS identification audit for labeled circle graphs was made by neither manuscript, nor by the write.
\item\label{cgr:q:moments} \emph{Unbounded automorphism observables} (Part~I, Proposition~\ref{cgr:ld:prop:aut} and Section~\ref{cgr:ld:sec:quantitative}; Part~II, Section~\ref{cgr:fc:sec:weighted}). Both Parts control only bounded functions of $|\Aut(G)|$ and say so. The limiting law $2^P$, $P\sim\Poi(1)$, has $\EE\,2^{tP}=\exp(2^t-1)<\infty$ for every $t$; whether $\EE|\Aut(G)|^t$ converges to it is open. What is missing is uniform integrability: control of the rare graphs with large automorphism groups (symmetric cores, large splits), which the bounded-weight arguments discard without bounding their weight.
\item\label{cgr:q:data} \emph{Exact counts beyond order~13} (write). The exact data stop at $g_{13}$ (Johnston, 2020) and the derived $c_{13}$. At $n=13$ the counts are still $10.5\%$ ($c_{13}$) and $11.3\%$ ($g_{13}$) below the first-corrected approximations (dated note in Section~\ref{cgr:ld:sec:checks}), so the expansions cannot yet be tested numerically, nor $\beta$ estimated. Part~II's Theorem~\ref{cgr:fc:thm:finite-transfer} expresses $c_n$ through asymmetric prime counts and rooted branch counts up to an error; an exact version, counting all prime centres with their port actions, would give a recurrence from the counts of prime circle graphs. Neither manuscript states one.
\end{enumerate}
\begin{thebibliography}{9}
\bibitem{BBFGP24}
F. Bassino, M. Bouvel, V. F\'eray, L. Gerin and A. Pierrot,
\emph{Dense and nondense limits for uniform random intersection
graphs}, arXiv:2402.06394v1, 9 February 2024.
\href{https://arxiv.org/abs/2402.06394v1}{Versioned arXiv record}.
Part~I: Sections 5.2--5.3 and Appendix A.2 are the primary inputs.
The 2026 journal metadata is identified in the text; its final
text was not inspected. Part~II: Sections 5.2--5.3 and Appendix A.2 are the inspected primary
circle-graph inputs; the final 2026 journal text was not inspected. \wtag\ Journal version: \emph{Electronic Journal of Probability} \textbf{31} (2026), article 110, \href{https://doi.org/10.1214/26-EJP1570}{doi:10.1214/26-EJP1570}; not read.
\bibitem{DP09}
L. E. Danielsen and M. G. Parker,
\emph{Interlace Polynomials: Enumeration, Unimodality, and
Connections to Codes}, arXiv:0804.2576v2, 1 October 2009,
Section 3, Table 3.
\href{https://arxiv.org/abs/0804.2576v2}{Versioned arXiv record}.
\bibitem{GP12}
E. Gioan and C. Paul,
\emph{Split decomposition and graph-labelled trees:
Characterizations and fully dynamic algorithms for totally
decomposable graphs}, Discrete Applied Mathematics
\textbf{160} (2012), 708--733.
\href{https://doi.org/10.1016/j.dam.2011.05.007}
{doi:10.1016/j.dam.2011.05.007}.
Section 2.2, especially Theorem 2.9 and Corollary 2.10. \wtag\ Also arXiv:0810.1823v2.
\bibitem{KZ15}
P. Klav\'ik and P. Zeman,
\emph{Automorphism Groups of Geometrically Represented Graphs},
32nd International Symposium on Theoretical Aspects of Computer
Science (STACS 2015), LIPIcs \textbf{30}, 540--553, 2015.
\href{https://doi.org/10.4230/LIPIcs.STACS.2015.540}
{doi:10.4230/LIPIcs.STACS.2015.540}.
Part~II: Section 3, circle graphs, canonical split trees and prime
automorphisms. \wtag\ Also arXiv:1407.2136.
\bibitem{OEISc}
OEIS Foundation Inc., \emph{The On-Line Encyclopedia of Integer
Sequences}, A156808, number of nonisomorphic connected circle
graphs of order $n$; inspected 3 October 2026 (Part~I).
\url{https://oeis.org/A156808}.
\bibitem{OEISg}
OEIS Foundation Inc., \emph{The On-Line Encyclopedia of Integer
Sequences}, A156809, number of nonisomorphic circle graphs of
order $n$; inspected 3 October 2026 (Part~I).
\url{https://oeis.org/A156809}.
\bibitem{TSvol}
\wtag\ \emph{Transseries and Inversion}, the transseries volume of this repository,
\file{Analysis/Transseries/docs/series-and-transseries/Transseries_And_Inversion/transseries_and_inversion.tex};
\texttt{p0:prop:factorial-core}, \texttt{p0:thm:lambert-core}, \texttt{p0:thm:staircase}.
\end{thebibliography}
\end{document}
